\documentclass[../main.tex]{subfiles}
\begin{document}
\newcommand{\lxOUT}{\mathrm{OUT}}\newcommand{\lxIN}{\mathrm{IN}}
\newcommand{\lxENTRY}{\mathrm{ENTRY}}\newcommand{\lxgen}{\mathrm{gen}}
\newcommand{\lxkill}{\mathrm{kill}}\newcommand{\lxpred}{\mathrm{pred}}

\chapter{Compiler Background - I}
\lessoninfo{
  video={M10, lessons 1--5 (GPU Compiler Flow; PTX; IR and Basic Block; Introduction to Data Flow Analysis; Example of Reaching Definitions)},
  sources={NVCC documentation \cite{nvidia2026nvcc}; PTX ISA 9.4 \cite{nvidia2026ptx}; Compiling CUDA with clang \cite{llvm2026compilecuda}; NVVM IR specification \cite{nvidia2026nvvm}; Lattner and Adve \cite{lattner2004llvm} (LLVM); LLVM Language Reference \cite{llvm2026langref}},
  textbook={K12 \cite{kim2012gpgpu} §1.3.7; CPG \cite{nvidia2024cudaguide} §3.1 Compilation with NVCC; Aho et al.\ \cite{aho2006compilers} §8.4, §9.2},
  keywords={compiler, front end, back end, intermediate representation, abstract syntax tree, nvcc, LLVM, NVPTX, PTX, SASS, virtual architecture, fat binary, JIT compilation, guard predicate, three-address code, leader, control flow graph, data-flow analysis, transfer function, reaching definitions, gen and kill sets, fixed point},
}

The previous lessons treated a CUDA kernel as something the GPU simply
runs; between the source file and the SM stands a compiler, and several
mechanisms of earlier lessons were in fact the compiler's work: the
predicated code of Lesson~6, the registers each thread receives in
Lesson~3, the PTX fragments of Lesson~4.  This lesson and the next give
the compiler background needed to reason about such code.  The first
half follows a CUDA program through the compilation flow, from source
code over an intermediate representation and PTX to the machine code
that the GPU executes, and looks at PTX itself.  The second half
introduces the structures on which optimisations work: three-address
code, basic blocks and control-flow graphs, and data-flow analysis,
worked out for the analysis called reaching definitions.  Lesson~11
continues with liveness, SSA form, classic optimisations and the
divergence analysis that connects all of this back to warps.

\section{GPU compiler flow}
\label{sec:l10-flow}

A \term{compiler}[コンパイラ] translates a program written in one
language, the source program, into an equivalent program in another
language, usually the instruction set of a machine.\source{M10, GPU
compiler flow; CPG §3.1 Compilation with NVCC; K12 §1.3.7;
\cite{nvidia2026nvcc}.}  It is organised as a pipeline of three parts
(\cref{fig:l10-compile-flow}).

\begin{figure}[tb]
  \centering
  % The classic compiler pipeline (Lesson 10).  A source program passes
% through the front end, which produces IR; the optimizer transforms IR
% into better IR; the back end generates code for the target machine.
% Redraws the outline figure m10l1_compileflow.
% Usage: % The classic compiler pipeline (Lesson 10).  A source program passes
% through the front end, which produces IR; the optimizer transforms IR
% into better IR; the back end generates code for the target machine.
% Redraws the outline figure m10l1_compileflow.
% Usage: % The classic compiler pipeline (Lesson 10).  A source program passes
% through the front end, which produces IR; the optimizer transforms IR
% into better IR; the back end generates code for the target machine.
% Redraws the outline figure m10l1_compileflow.
% Usage: \input{figures/l10-compile-flow}
% Width: body (120 mm); the drawing is about 114 mm wide.
\begin{tikzpicture}[x=1mm, y=1mm, font=\scriptsize\sffamily,
    >={Stealth[length=1.6mm]},
    st/.style={draw, fill=ln-accent!20, minimum width=20mm, minimum height=9mm,
               align=center, inner sep=1pt},
    io/.style={align=center, inner sep=1pt, text width=17mm}]
  \node[io] (src) at (8,0) {\iflangja{ソース\\プログラム}{source\\program}};
  \node[st] (fe) at (32,0) {\iflangja{フロント\\エンド}{front end}};
  \node[st] (op) at (61,0) {\iflangja{オプティマイザ}{optimizer}};
  \node[st] (be) at (90,0) {\iflangja{バックエンド\\（コード生成）}{back end\\(code generation)}};
  \node[io] (tgt) at (111,0) {\iflangja{ターゲット\\プログラム}{target\\program}};
  \draw[->] (src) -- (fe);
  \draw[->] (fe) -- node[above, inner sep=1.5pt] {IR} (op);
  \draw[->] (op) -- node[above, inner sep=1.5pt] {IR} (be);
  \draw[->] (be) -- (tgt);
  \draw[->] (op.north east) ++(-4,0) arc[start angle=-30, end angle=210, radius=4.6]
    node[pos=0.5, above, inner sep=1pt] {\iflangja{IR→IR}{IR$\to$IR}};
  \node[text=ln-muted, align=center] at (32,-9)
    {\iflangja{言語ごと}{one per language}};
  \node[text=ln-muted, align=center] at (61,-9)
    {\iflangja{共通}{shared}};
  \node[text=ln-muted, align=center] at (90,-9)
    {\iflangja{ターゲットごと}{one per target}};
\end{tikzpicture}

% Width: body (120 mm); the drawing is about 114 mm wide.
\begin{tikzpicture}[x=1mm, y=1mm, font=\scriptsize\sffamily,
    >={Stealth[length=1.6mm]},
    st/.style={draw, fill=ln-accent!20, minimum width=20mm, minimum height=9mm,
               align=center, inner sep=1pt},
    io/.style={align=center, inner sep=1pt, text width=17mm}]
  \node[io] (src) at (8,0) {\iflangja{ソース\\プログラム}{source\\program}};
  \node[st] (fe) at (32,0) {\iflangja{フロント\\エンド}{front end}};
  \node[st] (op) at (61,0) {\iflangja{オプティマイザ}{optimizer}};
  \node[st] (be) at (90,0) {\iflangja{バックエンド\\（コード生成）}{back end\\(code generation)}};
  \node[io] (tgt) at (111,0) {\iflangja{ターゲット\\プログラム}{target\\program}};
  \draw[->] (src) -- (fe);
  \draw[->] (fe) -- node[above, inner sep=1.5pt] {IR} (op);
  \draw[->] (op) -- node[above, inner sep=1.5pt] {IR} (be);
  \draw[->] (be) -- (tgt);
  \draw[->] (op.north east) ++(-4,0) arc[start angle=-30, end angle=210, radius=4.6]
    node[pos=0.5, above, inner sep=1pt] {\iflangja{IR→IR}{IR$\to$IR}};
  \node[text=ln-muted, align=center] at (32,-9)
    {\iflangja{言語ごと}{one per language}};
  \node[text=ln-muted, align=center] at (61,-9)
    {\iflangja{共通}{shared}};
  \node[text=ln-muted, align=center] at (90,-9)
    {\iflangja{ターゲットごと}{one per target}};
\end{tikzpicture}

% Width: body (120 mm); the drawing is about 114 mm wide.
\begin{tikzpicture}[x=1mm, y=1mm, font=\scriptsize\sffamily,
    >={Stealth[length=1.6mm]},
    st/.style={draw, fill=ln-accent!20, minimum width=20mm, minimum height=9mm,
               align=center, inner sep=1pt},
    io/.style={align=center, inner sep=1pt, text width=17mm}]
  \node[io] (src) at (8,0) {\iflangja{ソース\\プログラム}{source\\program}};
  \node[st] (fe) at (32,0) {\iflangja{フロント\\エンド}{front end}};
  \node[st] (op) at (61,0) {\iflangja{オプティマイザ}{optimizer}};
  \node[st] (be) at (90,0) {\iflangja{バックエンド\\（コード生成）}{back end\\(code generation)}};
  \node[io] (tgt) at (111,0) {\iflangja{ターゲット\\プログラム}{target\\program}};
  \draw[->] (src) -- (fe);
  \draw[->] (fe) -- node[above, inner sep=1.5pt] {IR} (op);
  \draw[->] (op) -- node[above, inner sep=1.5pt] {IR} (be);
  \draw[->] (be) -- (tgt);
  \draw[->] (op.north east) ++(-4,0) arc[start angle=-30, end angle=210, radius=4.6]
    node[pos=0.5, above, inner sep=1pt] {\iflangja{IR→IR}{IR$\to$IR}};
  \node[text=ln-muted, align=center] at (32,-9)
    {\iflangja{言語ごと}{one per language}};
  \node[text=ln-muted, align=center] at (61,-9)
    {\iflangja{共通}{shared}};
  \node[text=ln-muted, align=center] at (90,-9)
    {\iflangja{ターゲットごと}{one per target}};
\end{tikzpicture}

  \caption{The classic compiler pipeline.  The front end turns the
    source program into IR, the optimizer rewrites IR into better IR,
    and the back end generates code for the target machine.}
  \label{fig:l10-compile-flow}
\end{figure}

\begin{definition}[Front end, IR, back end]
  The \term{front end}[フロントエンド] reads the source program, checks
  its syntax and meaning, and translates it into an
  \term{intermediate representation (IR)}[中間表現], a program form
  that belongs neither to the source language nor to the
  target machine.\index{IR|see{intermediate representation (IR)}}%
  \margin{The optimizer is also called the \emph{middle end}; in
  LLVM it is a sequence of passes, each an IR-to-IR transformation
  \cite{lattner2004llvm}.}  The
  \emph{optimizer} transforms IR into IR that computes the same results
  faster or with fewer instructions.  The
  \term{back end}[バックエンド] performs
  \term{code generation}[コード生成]: it translates IR into
  instructions of the target machine and applies the optimisations that
  depend on that machine.
\end{definition}

The IR is what makes the split useful.  One front end per language and
one back end per target share a single optimizer, so that adding a
language or a machine adds one component rather than one complete
compiler per pair.\margin{With $L$ languages and $T$ targets, the IR
split needs $L+T$ front and back ends instead of $L\cdot T$
complete compilers.}  The IR is also a good level for optimisation:
it is simpler than the source language and still independent of any
machine, and it can carry metadata, such as types and source
positions, from the front end to the later stages.

\subsection{The flow for a GPU program}

A CUDA source file mixes host code and device code (Lesson~3).  A GPU
compiler therefore first separates the two and sends each through a
pipeline of its own: the device code is compiled for the GPU, the host
code with an ordinary C++ compiler, and the result for the GPU is
packed into the host program, which loads it when the program
runs.\margin{CPG §3.1.1.1 Offline Compilation \cite{nvidia2024cudaguide}:
nvcc separates device code from host code, compiles the device code to
PTX or binary form, and replaces the kernel launch syntax with calls
into the CUDA runtime.}

NVIDIA's compiler driver \term{nvcc}[nvcc] implements this flow as the
\emph{CUDA compilation trajectory} of \cref{fig:l10-cuda-flow}
\cite[§3]{nvidia2026nvcc}.  The input file is preprocessed twice.  On
the device side the front end and optimizer \texttt{cicc} produces a
\texttt{.ptx} file, and the assembler \texttt{ptxas} translates it into
a \texttt{.cubin} file of machine code; this pair of steps is repeated
for every architecture requested.  The tool \texttt{fatbinary} packs
the PTX and \texttt{.cubin} results into one image.  On the host side
\texttt{cudafe++} separates the host code and rewrites the CUDA
extensions into standard C++, and the host C++ compiler compiles that
code with the device image embedded.  The linker finally produces the
executable.\margin{The tool names come from the trajectory figure of
the NVCC documentation \cite[§3]{nvidia2026nvcc}; \texttt{nvcc
-{}-dryrun} lists the commands it would run, and \texttt{-{}-keep}
keeps the intermediate \texttt{.ptx} and \texttt{.cubin} files
(§4.2.5.9, §4.2.5.11).}  With separate compilation of device code, a
device linker \texttt{nvlink} first links the \texttt{.cubin} files of
all translation units.

\begin{figure}[tb]
  \centering
  % The CUDA compilation trajectory of nvcc, simplified (Lesson 10).  A
% .cu file is preprocessed twice.  Device path: cicc (front end,
% optimizer, PTX code generation) -> .ptx -> ptxas -> .cubin (SASS),
% repeated per architecture; fatbinary packs PTX and cubin into a fat
% binary.  Host path: cudafe++ separates host code, the host C++ compiler
% compiles it with the fat binary embedded; the linker builds the
% executable.  The device linker nvlink (separate compilation) is omitted.
% Redraws the outline figures m10l1_compileflowgpu and
% m10l1_compileflowgpunvidia (simplified; not NVIDIA's figure).
% Usage: % The CUDA compilation trajectory of nvcc, simplified (Lesson 10).  A
% .cu file is preprocessed twice.  Device path: cicc (front end,
% optimizer, PTX code generation) -> .ptx -> ptxas -> .cubin (SASS),
% repeated per architecture; fatbinary packs PTX and cubin into a fat
% binary.  Host path: cudafe++ separates host code, the host C++ compiler
% compiles it with the fat binary embedded; the linker builds the
% executable.  The device linker nvlink (separate compilation) is omitted.
% Redraws the outline figures m10l1_compileflowgpu and
% m10l1_compileflowgpunvidia (simplified; not NVIDIA's figure).
% Usage: % The CUDA compilation trajectory of nvcc, simplified (Lesson 10).  A
% .cu file is preprocessed twice.  Device path: cicc (front end,
% optimizer, PTX code generation) -> .ptx -> ptxas -> .cubin (SASS),
% repeated per architecture; fatbinary packs PTX and cubin into a fat
% binary.  Host path: cudafe++ separates host code, the host C++ compiler
% compiles it with the fat binary embedded; the linker builds the
% executable.  The device linker nvlink (separate compilation) is omitted.
% Redraws the outline figures m10l1_compileflowgpu and
% m10l1_compileflowgpunvidia (simplified; not NVIDIA's figure).
% Usage: \input{figures/l10-cuda-flow}
% Width: body (120 mm); the drawing is about 116 mm wide.
\begin{tikzpicture}[x=1mm, y=1mm, font=\scriptsize\sffamily,
    >={Stealth[length=1.5mm]},
    tool/.style={draw, fill=ln-accent!20, minimum height=8mm, align=center,
                 inner sep=1.5pt},
    file/.style={draw, fill=ln-box, rounded corners=1.5pt, minimum height=6mm,
                 align=center, inner sep=1.5pt, font=\scriptsize\ttfamily}]
  % source
  \node[file, minimum height=10mm] (cu) at (6,20) {x.cu};
  % device row
  \node[tool] (dpp) at (22,36) {\iflangja{プリ\\プロセッサ}{pre-\\processor}};
  \node[tool, fill=ln-source!25] (cicc) at (44,36) {cicc\\[-1pt]
    {\scriptsize\sffamily\iflangja{FE＋最適化＋PTX生成}{FE + opt.\ + PTX gen.}}};
  \node[file] (ptx) at (65,36) {x.ptx};
  \node[tool, fill=ln-source!25] (ptxas) at (81,36) {ptxas};
  \node[file] (cubin) at (100,36) {x.cubin\\[-1pt]{\scriptsize\sffamily (SASS)}};
  % fatbinary
  \node[tool] (fatb) at (100,20) {fatbinary};
  % host row
  \node[tool] (hpp) at (22,4) {\iflangja{プリ\\プロセッサ}{pre-\\processor}};
  \node[tool] (fe) at (44,4) {cudafe++};
  \node[tool] (cxx) at (71,4) {\iflangja{ホスト C++\\コンパイラ}{host C++\\compiler}};
  \node[file] (obj) at (91,4) {x.o};
  \node[tool] (ld) at (108,4) {\iflangja{リンカ}{linker}};
  \draw[->] (cu.north) |- (dpp.west);
  \draw[->] (cu.south) |- (hpp.west);
  \draw[->] (dpp) -- (cicc);
  \draw[->] (cicc) -- (ptx);
  \draw[->] (ptx) -- (ptxas);
  \draw[->] (ptxas) -- (cubin);
  \draw[->] (cubin) -- (fatb);
  \draw[->] (ptx.south) |- (fatb.west);
  \draw[->] (hpp) -- (fe);
  \draw[->] (fe) -- (cxx);
  \draw[->] (cxx) -- (obj);
  \draw[->] (obj) -- (ld);
  \draw[->, dashed] (fatb.south) -- ++(0,-4) -| node[pos=0.25, above, inner sep=1pt]
    {\iflangja{ファットバイナリを埋め込む}{fat binary embedded}} (cxx.north);
  \draw[->] (ld.south) -- ++(0,-6) node[below, inner sep=1pt]
    {\iflangja{実行ファイル}{executable}};
  % lane labels
  \node[anchor=south west, text=ln-muted] at (13,41.5)
    {\iflangja{デバイスコード（アーキテクチャごとに繰り返す）}{device code (repeated per architecture)}};
  \node[anchor=north west, text=ln-muted] at (13,-1.5)
    {\iflangja{ホストコード}{host code}};
\end{tikzpicture}

% Width: body (120 mm); the drawing is about 116 mm wide.
\begin{tikzpicture}[x=1mm, y=1mm, font=\scriptsize\sffamily,
    >={Stealth[length=1.5mm]},
    tool/.style={draw, fill=ln-accent!20, minimum height=8mm, align=center,
                 inner sep=1.5pt},
    file/.style={draw, fill=ln-box, rounded corners=1.5pt, minimum height=6mm,
                 align=center, inner sep=1.5pt, font=\scriptsize\ttfamily}]
  % source
  \node[file, minimum height=10mm] (cu) at (6,20) {x.cu};
  % device row
  \node[tool] (dpp) at (22,36) {\iflangja{プリ\\プロセッサ}{pre-\\processor}};
  \node[tool, fill=ln-source!25] (cicc) at (44,36) {cicc\\[-1pt]
    {\scriptsize\sffamily\iflangja{FE＋最適化＋PTX生成}{FE + opt.\ + PTX gen.}}};
  \node[file] (ptx) at (65,36) {x.ptx};
  \node[tool, fill=ln-source!25] (ptxas) at (81,36) {ptxas};
  \node[file] (cubin) at (100,36) {x.cubin\\[-1pt]{\scriptsize\sffamily (SASS)}};
  % fatbinary
  \node[tool] (fatb) at (100,20) {fatbinary};
  % host row
  \node[tool] (hpp) at (22,4) {\iflangja{プリ\\プロセッサ}{pre-\\processor}};
  \node[tool] (fe) at (44,4) {cudafe++};
  \node[tool] (cxx) at (71,4) {\iflangja{ホスト C++\\コンパイラ}{host C++\\compiler}};
  \node[file] (obj) at (91,4) {x.o};
  \node[tool] (ld) at (108,4) {\iflangja{リンカ}{linker}};
  \draw[->] (cu.north) |- (dpp.west);
  \draw[->] (cu.south) |- (hpp.west);
  \draw[->] (dpp) -- (cicc);
  \draw[->] (cicc) -- (ptx);
  \draw[->] (ptx) -- (ptxas);
  \draw[->] (ptxas) -- (cubin);
  \draw[->] (cubin) -- (fatb);
  \draw[->] (ptx.south) |- (fatb.west);
  \draw[->] (hpp) -- (fe);
  \draw[->] (fe) -- (cxx);
  \draw[->] (cxx) -- (obj);
  \draw[->] (obj) -- (ld);
  \draw[->, dashed] (fatb.south) -- ++(0,-4) -| node[pos=0.25, above, inner sep=1pt]
    {\iflangja{ファットバイナリを埋め込む}{fat binary embedded}} (cxx.north);
  \draw[->] (ld.south) -- ++(0,-6) node[below, inner sep=1pt]
    {\iflangja{実行ファイル}{executable}};
  % lane labels
  \node[anchor=south west, text=ln-muted] at (13,41.5)
    {\iflangja{デバイスコード（アーキテクチャごとに繰り返す）}{device code (repeated per architecture)}};
  \node[anchor=north west, text=ln-muted] at (13,-1.5)
    {\iflangja{ホストコード}{host code}};
\end{tikzpicture}

% Width: body (120 mm); the drawing is about 116 mm wide.
\begin{tikzpicture}[x=1mm, y=1mm, font=\scriptsize\sffamily,
    >={Stealth[length=1.5mm]},
    tool/.style={draw, fill=ln-accent!20, minimum height=8mm, align=center,
                 inner sep=1.5pt},
    file/.style={draw, fill=ln-box, rounded corners=1.5pt, minimum height=6mm,
                 align=center, inner sep=1.5pt, font=\scriptsize\ttfamily}]
  % source
  \node[file, minimum height=10mm] (cu) at (6,20) {x.cu};
  % device row
  \node[tool] (dpp) at (22,36) {\iflangja{プリ\\プロセッサ}{pre-\\processor}};
  \node[tool, fill=ln-source!25] (cicc) at (44,36) {cicc\\[-1pt]
    {\scriptsize\sffamily\iflangja{FE＋最適化＋PTX生成}{FE + opt.\ + PTX gen.}}};
  \node[file] (ptx) at (65,36) {x.ptx};
  \node[tool, fill=ln-source!25] (ptxas) at (81,36) {ptxas};
  \node[file] (cubin) at (100,36) {x.cubin\\[-1pt]{\scriptsize\sffamily (SASS)}};
  % fatbinary
  \node[tool] (fatb) at (100,20) {fatbinary};
  % host row
  \node[tool] (hpp) at (22,4) {\iflangja{プリ\\プロセッサ}{pre-\\processor}};
  \node[tool] (fe) at (44,4) {cudafe++};
  \node[tool] (cxx) at (71,4) {\iflangja{ホスト C++\\コンパイラ}{host C++\\compiler}};
  \node[file] (obj) at (91,4) {x.o};
  \node[tool] (ld) at (108,4) {\iflangja{リンカ}{linker}};
  \draw[->] (cu.north) |- (dpp.west);
  \draw[->] (cu.south) |- (hpp.west);
  \draw[->] (dpp) -- (cicc);
  \draw[->] (cicc) -- (ptx);
  \draw[->] (ptx) -- (ptxas);
  \draw[->] (ptxas) -- (cubin);
  \draw[->] (cubin) -- (fatb);
  \draw[->] (ptx.south) |- (fatb.west);
  \draw[->] (hpp) -- (fe);
  \draw[->] (fe) -- (cxx);
  \draw[->] (cxx) -- (obj);
  \draw[->] (obj) -- (ld);
  \draw[->, dashed] (fatb.south) -- ++(0,-4) -| node[pos=0.25, above, inner sep=1pt]
    {\iflangja{ファットバイナリを埋め込む}{fat binary embedded}} (cxx.north);
  \draw[->] (ld.south) -- ++(0,-6) node[below, inner sep=1pt]
    {\iflangja{実行ファイル}{executable}};
  % lane labels
  \node[anchor=south west, text=ln-muted] at (13,41.5)
    {\iflangja{デバイスコード（アーキテクチャごとに繰り返す）}{device code (repeated per architecture)}};
  \node[anchor=north west, text=ln-muted] at (13,-1.5)
    {\iflangja{ホストコード}{host code}};
\end{tikzpicture}

  \caption{The nvcc compilation trajectory, simplified.  Device code
    passes cicc and ptxas once per architecture; fatbinary packs the
    PTX and the machine code (cubin) into a fat binary that the host
    object embeds.  The device linker nvlink is omitted.}
  \label{fig:l10-cuda-flow}
\end{figure}

\begin{definition}[Virtual architecture]
  A \term{virtual architecture}[仮想アーキテクチャ] is the GPU that PTX
  code is written for: a set of capabilities that a program may use,
  named \texttt{compute\_XY} after a compute capability~X.Y.  An actual
  GPU, the target of \texttt{ptxas}, is a \emph{real architecture},
  named \texttt{sm\_XY}.
\end{definition}

Every nvcc command therefore names two architectures, a virtual one for
the first stage and a real one for the second, and the real one must
implement the virtual one \cite[§5.4]{nvidia2026nvcc}.\tip{Choose the
virtual architecture as low as possible, so that the PTX runs on as
many GPUs as possible, and the real architecture as high as the target
GPUs allow (NVCC §5.4 \cite{nvidia2026nvcc}).}  The virtual
architecture states which features the program needs; it does not fix
the GPU it will run on.

\subsection{The open-source toolchain: Clang and LLVM}

The same structure appears in the open-source CUDA compiler.

\begin{definition}[LLVM, Clang, NVPTX]
  \term{LLVM}[LLVM] is a compiler infrastructure built around one typed
  IR, the LLVM IR, with a library of optimisations on it and back ends
  for many targets \cite{lattner2004llvm}.
  \term{Clang}[Clang] is the LLVM front end for the C family of
  languages, including CUDA.  \term{NVPTX}[NVPTX] is the LLVM back end
  that generates PTX.
\end{definition}

A splitter separates host and device code.  The device code passes
through Clang, the IR optimisations and the NVPTX code generator and
becomes PTX assembly; the host compilation then embeds the device code
into a fat binary.\sidenote{Clang does not split the file before
parsing it: it uses \enquote{merged parsing}, in which all host and
device code is present in both compilation steps and must be
semantically correct in both.  Unlike nvcc, it always also runs
\texttt{ptxas} and embeds machine code for each requested GPU
(\texttt{-{}-offload-arch=sm\_80}).  By default it does not embed
PTX; \texttt{-{}-cuda-include-ptx=sm\_80} (or \texttt{all}) adds it
for forward compatibility \cite{llvm2026compilecuda}.}  NVIDIA's
own device compiler is built on the same technology: its front end
emits \emph{NVVM IR}, an IR \enquote{based on the LLVM IR}, and the
NVVM compiler, itself based on LLVM, generates PTX from it
\cite{nvidia2026nvvm}.

The front end parses the source into a tree.

\begin{definition}[Abstract syntax tree]
  An \term{abstract syntax tree (AST)}[抽象構文木] represents the
  structure of a program as a tree: each inner node is an operation or
  statement (a sequence, a loop, a branch, an assignment, an
  arithmetic operator) and its children are its operands or
  sub-statements; the leaves are variables and
  constants.\index{AST|see{abstract syntax tree (AST)}}
  Punctuation such as parentheses and semicolons does not appear: the
  tree shape already expresses it.
\end{definition}

\begin{example}[The AST of a loop]
  \label{ex:l10-ast}
  The loop below keeps the running maximum of an array.
\begin{lstlisting}[language=C++, numbers=none]
while (i < n) {
  if (x[i] > m) m = x[i];
  i = i + 1;
}
return m;
\end{lstlisting}
  Its AST (\cref{fig:l10-ast}) has a sequence at the root with two
  children, the \texttt{while} loop and the \texttt{return}.  The loop
  node has the condition \texttt{i < n} and the body, itself a sequence
  of the \texttt{if} statement and the increment.  Each array access
  \texttt{x[i]} is an indexing node with the array and the index as
  children.
\end{example}

\begin{figure}[tb]
  \centering
  % Abstract syntax tree of a running-maximum loop (Lesson 10).  The tree of
%   while (i < n) { if (x[i] > m) m = x[i]; i = i + 1; }  return m;
% Operations (sequence, while, if, assignment, comparison, addition,
% indexing) are blue, variables green, constants red.
% Original (replaces the outline figure m10l1_clang, which shows GCD).
% Usage: % Abstract syntax tree of a running-maximum loop (Lesson 10).  The tree of
%   while (i < n) { if (x[i] > m) m = x[i]; i = i + 1; }  return m;
% Operations (sequence, while, if, assignment, comparison, addition,
% indexing) are blue, variables green, constants red.
% Original (replaces the outline figure m10l1_clang, which shows GCD).
% Usage: % Abstract syntax tree of a running-maximum loop (Lesson 10).  The tree of
%   while (i < n) { if (x[i] > m) m = x[i]; i = i + 1; }  return m;
% Operations (sequence, while, if, assignment, comparison, addition,
% indexing) are blue, variables green, constants red.
% Original (replaces the outline figure m10l1_clang, which shows GCD).
% Usage: \input{figures/l10-ast}
% Width: body (120 mm); the drawing is about 116 mm wide.
\begin{tikzpicture}[x=1mm, y=1mm, font=\scriptsize\sffamily,
    op/.style={draw, fill=ln-accent!25, rounded corners=1.5pt, minimum height=5mm,
               minimum width=7mm, inner sep=1.5pt},
    var/.style={draw, fill=ln-tip!30, minimum height=5mm, minimum width=5mm, inner sep=1.5pt},
    cst/.style={draw, fill=ln-caution!25, minimum height=5mm, minimum width=5mm, inner sep=1.5pt}]
  \node[op] (seq) at (55,50) {\iflangja{文の並び}{sequence}};
  \node[op] (wh) at (40,40) {while};
  \node[op] (ret) at (92,40) {return};
  \node[var] (rm) at (92,30) {\texttt{m}};
  \node[op] (lt) at (12,30) {\texttt{<}};
  \node[op] (body) at (55,30) {\iflangja{文の並び}{sequence}};
  \node[var] (li) at (5,20) {\texttt{i}};
  \node[var] (ln) at (19,20) {\texttt{n}};
  \node[op] (if) at (38,20) {if};
  \node[op] (as2) at (82,20) {\texttt{:=}};
  \node[op] (gt) at (28,10) {\texttt{>}};
  \node[op] (as1) at (54,10) {\texttt{:=}};
  \node[var] (ai) at (73,10) {\texttt{i}};
  \node[op] (pl) at (91,10) {\texttt{+}};
  \node[op] (ix1) at (19,0) {\texttt{[\,]}};
  \node[var] (gm) at (34,0) {\texttt{m}};
  \node[var] (am) at (46,0) {\texttt{m}};
  \node[op] (ix2) at (60,0) {\texttt{[\,]}};
  \node[var] (pi) at (85,0) {\texttt{i}};
  \node[cst] (p1) at (97,0) {\texttt{1}};
  \node[var] (x1) at (13,-10) {\texttt{x}};
  \node[var] (i1) at (25,-10) {\texttt{i}};
  \node[var] (x2) at (54,-10) {\texttt{x}};
  \node[var] (i2) at (66,-10) {\texttt{i}};
  \foreach \a/\b in {seq/wh, seq/ret, ret/rm, wh/lt, wh/body, lt/li, lt/ln,
                     body/if, body/as2, if/gt, if/as1, gt/ix1, gt/gm, as1/am, as1/ix2,
                     as2/ai, as2/pl, pl/pi, pl/p1, ix1/x1, ix1/i1, ix2/x2, ix2/i2}
    \draw (\a.south) -- (\b.north);
  \node[text=ln-muted, anchor=east, inner sep=1pt] at (25,36) {\iflangja{条件}{condition}};
  \node[text=ln-muted, anchor=west, inner sep=1pt] at (48,36) {\iflangja{本体}{body}};
  % legend
  \node[op, anchor=west] at (97,56) {};
  \node[anchor=west] at (104,56) {\iflangja{演算・文}{operation}};
  \node[var, anchor=west] at (97,50) {};
  \node[anchor=west] at (104,50) {\iflangja{変数}{variable}};
  \node[cst, anchor=west] at (97,44) {};
  \node[anchor=west] at (104,44) {\iflangja{定数}{constant}};
\end{tikzpicture}

% Width: body (120 mm); the drawing is about 116 mm wide.
\begin{tikzpicture}[x=1mm, y=1mm, font=\scriptsize\sffamily,
    op/.style={draw, fill=ln-accent!25, rounded corners=1.5pt, minimum height=5mm,
               minimum width=7mm, inner sep=1.5pt},
    var/.style={draw, fill=ln-tip!30, minimum height=5mm, minimum width=5mm, inner sep=1.5pt},
    cst/.style={draw, fill=ln-caution!25, minimum height=5mm, minimum width=5mm, inner sep=1.5pt}]
  \node[op] (seq) at (55,50) {\iflangja{文の並び}{sequence}};
  \node[op] (wh) at (40,40) {while};
  \node[op] (ret) at (92,40) {return};
  \node[var] (rm) at (92,30) {\texttt{m}};
  \node[op] (lt) at (12,30) {\texttt{<}};
  \node[op] (body) at (55,30) {\iflangja{文の並び}{sequence}};
  \node[var] (li) at (5,20) {\texttt{i}};
  \node[var] (ln) at (19,20) {\texttt{n}};
  \node[op] (if) at (38,20) {if};
  \node[op] (as2) at (82,20) {\texttt{:=}};
  \node[op] (gt) at (28,10) {\texttt{>}};
  \node[op] (as1) at (54,10) {\texttt{:=}};
  \node[var] (ai) at (73,10) {\texttt{i}};
  \node[op] (pl) at (91,10) {\texttt{+}};
  \node[op] (ix1) at (19,0) {\texttt{[\,]}};
  \node[var] (gm) at (34,0) {\texttt{m}};
  \node[var] (am) at (46,0) {\texttt{m}};
  \node[op] (ix2) at (60,0) {\texttt{[\,]}};
  \node[var] (pi) at (85,0) {\texttt{i}};
  \node[cst] (p1) at (97,0) {\texttt{1}};
  \node[var] (x1) at (13,-10) {\texttt{x}};
  \node[var] (i1) at (25,-10) {\texttt{i}};
  \node[var] (x2) at (54,-10) {\texttt{x}};
  \node[var] (i2) at (66,-10) {\texttt{i}};
  \foreach \a/\b in {seq/wh, seq/ret, ret/rm, wh/lt, wh/body, lt/li, lt/ln,
                     body/if, body/as2, if/gt, if/as1, gt/ix1, gt/gm, as1/am, as1/ix2,
                     as2/ai, as2/pl, pl/pi, pl/p1, ix1/x1, ix1/i1, ix2/x2, ix2/i2}
    \draw (\a.south) -- (\b.north);
  \node[text=ln-muted, anchor=east, inner sep=1pt] at (25,36) {\iflangja{条件}{condition}};
  \node[text=ln-muted, anchor=west, inner sep=1pt] at (48,36) {\iflangja{本体}{body}};
  % legend
  \node[op, anchor=west] at (97,56) {};
  \node[anchor=west] at (104,56) {\iflangja{演算・文}{operation}};
  \node[var, anchor=west] at (97,50) {};
  \node[anchor=west] at (104,50) {\iflangja{変数}{variable}};
  \node[cst, anchor=west] at (97,44) {};
  \node[anchor=west] at (104,44) {\iflangja{定数}{constant}};
\end{tikzpicture}

% Width: body (120 mm); the drawing is about 116 mm wide.
\begin{tikzpicture}[x=1mm, y=1mm, font=\scriptsize\sffamily,
    op/.style={draw, fill=ln-accent!25, rounded corners=1.5pt, minimum height=5mm,
               minimum width=7mm, inner sep=1.5pt},
    var/.style={draw, fill=ln-tip!30, minimum height=5mm, minimum width=5mm, inner sep=1.5pt},
    cst/.style={draw, fill=ln-caution!25, minimum height=5mm, minimum width=5mm, inner sep=1.5pt}]
  \node[op] (seq) at (55,50) {\iflangja{文の並び}{sequence}};
  \node[op] (wh) at (40,40) {while};
  \node[op] (ret) at (92,40) {return};
  \node[var] (rm) at (92,30) {\texttt{m}};
  \node[op] (lt) at (12,30) {\texttt{<}};
  \node[op] (body) at (55,30) {\iflangja{文の並び}{sequence}};
  \node[var] (li) at (5,20) {\texttt{i}};
  \node[var] (ln) at (19,20) {\texttt{n}};
  \node[op] (if) at (38,20) {if};
  \node[op] (as2) at (82,20) {\texttt{:=}};
  \node[op] (gt) at (28,10) {\texttt{>}};
  \node[op] (as1) at (54,10) {\texttt{:=}};
  \node[var] (ai) at (73,10) {\texttt{i}};
  \node[op] (pl) at (91,10) {\texttt{+}};
  \node[op] (ix1) at (19,0) {\texttt{[\,]}};
  \node[var] (gm) at (34,0) {\texttt{m}};
  \node[var] (am) at (46,0) {\texttt{m}};
  \node[op] (ix2) at (60,0) {\texttt{[\,]}};
  \node[var] (pi) at (85,0) {\texttt{i}};
  \node[cst] (p1) at (97,0) {\texttt{1}};
  \node[var] (x1) at (13,-10) {\texttt{x}};
  \node[var] (i1) at (25,-10) {\texttt{i}};
  \node[var] (x2) at (54,-10) {\texttt{x}};
  \node[var] (i2) at (66,-10) {\texttt{i}};
  \foreach \a/\b in {seq/wh, seq/ret, ret/rm, wh/lt, wh/body, lt/li, lt/ln,
                     body/if, body/as2, if/gt, if/as1, gt/ix1, gt/gm, as1/am, as1/ix2,
                     as2/ai, as2/pl, pl/pi, pl/p1, ix1/x1, ix1/i1, ix2/x2, ix2/i2}
    \draw (\a.south) -- (\b.north);
  \node[text=ln-muted, anchor=east, inner sep=1pt] at (25,36) {\iflangja{条件}{condition}};
  \node[text=ln-muted, anchor=west, inner sep=1pt] at (48,36) {\iflangja{本体}{body}};
  % legend
  \node[op, anchor=west] at (97,56) {};
  \node[anchor=west] at (104,56) {\iflangja{演算・文}{operation}};
  \node[var, anchor=west] at (97,50) {};
  \node[anchor=west] at (104,50) {\iflangja{変数}{variable}};
  \node[cst, anchor=west] at (97,44) {};
  \node[anchor=west] at (104,44) {\iflangja{定数}{constant}};
\end{tikzpicture}

  \caption{The abstract syntax tree of the loop in
    \cref{ex:l10-ast}.  Operations and statements are blue, variables
    green, constants red.}
  \label{fig:l10-ast}
\end{figure}

\tip{\texttt{clang -fsyntax-only -Xclang -ast-dump file.c} prints the
AST that Clang builds for a file.}  Before the parser runs, the
\term{preprocessor}[プリプロセッサ] of C and C++ performs pure text
substitution: it expands \texttt{\#include} files and macros and
evaluates conditional compilation.\margin{Since the preprocessor works
on text, a macro body such as \texttt{a + b} used in
\texttt{2 * M} becomes \texttt{2 * a + b}; macro bodies are
therefore written in parentheses.}  With the definition
\begin{lstlisting}[language=C++, numbers=none]
#define BLOCK_SIZE 256
int threads = BLOCK_SIZE;
\end{lstlisting}
the compiler proper sees only \texttt{int threads = 256;}.

\subsection{PTX, SASS and the fat binary}

The output of the device back end is not yet machine code.

\begin{definition}[PTX and SASS]
  \term{PTX (Parallel Thread Execution)}[PTX]\index{Parallel Thread Execution|see{PTX (Parallel Thread Execution)}} is NVIDIA's virtual ISA
  (Lesson~1): an architecture-independent instruction set whose
  programs are translated into machine code later.  Its registers are
  virtual, so PTX code has no register allocation.
  \term{SASS}[SASS] (shader assembly) is the machine code of one GPU
  architecture, with physical registers allocated.  SASS changes with
  every GPU generation; PTX stays compatible: code written for an older
  PTX version still compiles for newer GPUs, while new PTX versions add
  instructions.\margin{PTX ISA~8.0, the \texttt{.version} of
  \cref{ex:l10-ptx-kernel}, came with CUDA~12.0; the release history
  of PTX ISA~9.4 runs from PTX ISA~1.0 with CUDA~1.0
  \cite[§13]{nvidia2026ptx}.}
\end{definition}

The PTX ISA states the purpose directly: \enquote{PTX defines a virtual
machine and ISA for general purpose parallel thread execution}, and its
goals include \enquote{a stable ISA that spans multiple GPU
generations} \cite[§1.1--1.2]{nvidia2026ptx}.\margin{A PTX
register is a variable of a type, and a kernel may declare any number;
\texttt{ptxas} maps them onto at most 255 physical registers per thread
(CPG Table~21 \cite{nvidia2024cudaguide}) and spills the rest to local
memory \cite[§5.1.1]{nvidia2026ptx}.  Lesson~11 treats register
allocation.}

\begin{definition}[Fat binary]
  A \term{fat binary}[ファットバイナリ] is an executable that carries,
  next to the host code, several translations of the same device code:
  machine code for some real architectures and PTX for some virtual
  ones.  At run time the device part is handed to the driver, which
  selects the most suitable translation for the GPU present
  \cite[§5.6.2]{nvidia2026nvcc}.
\end{definition}

If the fat binary contains no machine code for the GPU present, the
driver compiles the PTX at run time.

\begin{definition}[JIT compilation]
  In \term{JIT compilation}[JITコンパイル] (just-in-time compilation)
  the device driver compiles PTX into machine code when an application
  loads it.  CPG §3.1.1.2 Just-in-Time Compilation states that it is
  \enquote{the only way for applications to run on devices that did not
  exist at the time the application was compiled}
  \cite{nvidia2024cudaguide}.
\end{definition}

JIT compilation costs start-up time.  The driver therefore keeps a copy
of each result in the \emph{compute cache}, which is invalidated when
the driver is upgraded, so that a new JIT compiler can improve old
applications (CPG §3.1.1.2).\margin{The compute cache lives in
\texttt{\textasciitilde/.nv/ComputeCache} on Linux and holds
\SI{1}{\gibi\byte} by default (\SI{256}{\mebi\byte} on embedded
platforms), at most \SI{4}{\gibi\byte}, set by
\texttt{CUDA\_CACHE\_MAXSIZE} (CPG §18 CUDA Environment Variables
\cite{nvidia2024cudaguide}).}  A fat binary that also carries machine
code for the expected GPUs avoids the delay altogether, and its PTX
keeps the program running on later GPUs:
\begin{lstlisting}[language=bash, numbers=none]
nvcc x.cu --gpu-architecture=compute_80 \
          --gpu-code=compute_80,sm_86,sm_89
\end{lstlisting}
builds machine code for compute capabilities 8.6 and~8.9 and keeps PTX
for compute capability~8.0, which a later GPU can JIT-compile
\cite[§5.6.2]{nvidia2026nvcc}.\caution{Machine code does not move
across major compute capabilities: a cubin for compute capability X.y
runs only on X.z with $z\ge y$ (CPG §3.1.2 Binary Compatibility
\cite{nvidia2024cudaguide}).}  PTX, in contrast, always compiles to an
equal or higher compute capability, but the result cannot use features
added after its virtual architecture: PTX for compute capability 6.0,
JIT-compiled on a 7.0 GPU, uses no tensor cores (CPG §3.1.3 PTX
Compatibility).

\section{PTX}
\label{sec:l10-ptx}

PTX is the lowest level at which CUDA code is normally read, so its
instruction format is worth knowing in detail.\source{M10, PTX; PTX ISA
§4.3, §9.2--9.3 \cite{nvidia2026ptx}; K12 §1.3.7.}

\subsection{Instruction format}

A PTX instruction consists of an \term{opcode}[オペコード], the name of
the operation, and a list of \term{operand}[オペランド]s.  The PTX ISA
puts it as follows: \enquote{PTX instructions generally have from zero
to four operands, plus an optional guard predicate appearing after an
\texttt{@} symbol to the left of the opcode}
\cite[§9.2]{nvidia2026ptx}.
\begin{lstlisting}[language=PTX, numbers=none]
@p  opcode;
@p  opcode a;
@p  opcode d, a;
@p  opcode d, a, b;
@p  opcode d, a, b, c;
\end{lstlisting}
For an instruction that produces a value, \texttt{d} is the
destination and \texttt{a}, \texttt{b}, \texttt{c} are sources; the
destination comes first.\margin{The opcode carries suffixes: a type
(\texttt{.s32} signed, \texttt{.u32} unsigned, \texttt{.f32} float,
\texttt{.b32} untyped bits, \texttt{.pred}) and, for memory
instructions, a state space (\texttt{.global}, \texttt{.shared},
\texttt{.local}, \texttt{.param}).}  Every arithmetic and memory
instruction names the type it operates on, as in \texttt{add.s32} or \texttt{ld.global.f32}.

\begin{definition}[Guard predicate, predicate register]
  A \term{predicate register}[プレディケートレジスタ] is a one-bit PTX register
  of type \texttt{.pred}, declared like any other register, e.g.\
  \texttt{.reg .pred p, q;}.  A \term{guard predicate}[ガードプレディケート]
  \texttt{@p} placed before an instruction makes it conditional: the
  instruction takes effect only in threads where \texttt{p} is true;
  \texttt{@!p} uses the negation.  An instruction without a guard is
  executed unconditionally.
\end{definition}

Predicate registers are virtual like all PTX registers; the
\texttt{.pred} type tells \texttt{ptxas} that one bit
suffices \cite[§9.3]{nvidia2026ptx}.\margin{Predicates are the
mechanism behind the if-conversion and predicated execution of
Lessons 5 and~6.}  They are set most commonly by \texttt{setp}, which
\enquote{compare[s] two numeric values with a relational operator}
\cite[§9.7.7.2]{nvidia2026ptx}.  \texttt{setp} can write two
destinations, separated by \texttt{|} (a few other instructions, such
as \texttt{shfl.sync}, use the same notation): the second receives the
negation of the comparison.

\begin{example}[Removing a branch with predicates]
  \label{ex:l10-predication}
  The statement \texttt{if (a < b) c = c + 1; else c = c - 1;} with
  \texttt{a}, \texttt{b}, \texttt{c} in \texttt{\%r1}, \texttt{\%r2},
  \texttt{\%r3} needs no branch:
\begin{lstlisting}[language=PTX, numbers=none]
     setp.lt.s32  %p1|%p2, %r1, %r2; // p1 = a<b, p2 = !p1
@%p1 add.s32      %r3, %r3, 1;       // then: c = c + 1
@%p2 sub.s32      %r3, %r3, 1;       // else: c = c - 1
\end{lstlisting}
  Every thread of a warp issues all three instructions, and each thread
  commits exactly one of the two updates.  The same guard on a branch
  instruction gives a real conditional jump:
\begin{lstlisting}[language=PTX, numbers=none]
     setp.lt.s32  %p1, %r1, %r2;
@!%p1 bra         ELSE;              // a >= b: jump
\end{lstlisting}
\end{example}

\tip{The comparison operators are \texttt{eq}, \texttt{ne},
\texttt{lt}, \texttt{le}, \texttt{gt}, \texttt{ge} for signed integers
and floats; unsigned integers use \texttt{eq}, \texttt{ne},
\texttt{lo}, \texttt{ls}, \texttt{hi}, \texttt{hs} (lower, lower or
same, higher, higher or same) \cite[§9.3.1]{nvidia2026ptx}.}  The predicated form trades the
risk of divergence (Lesson~6) for issuing both sides; for short bodies
like these the compiler prefers it.

\subsection{Statements}

A PTX program is a sequence of statements.  In the words of the PTX
ISA, \enquote{a PTX statement is either a directive or an instruction.
Statements begin with an optional label and end with a semicolon}
\cite[§4.3]{nvidia2026ptx}.

\begin{definition}[Directive, label]
  A \term{directive}[ディレクティブ] is a statement whose keyword begins
  with a dot and which declares rather than computes: the PTX version
  and target (\texttt{.version}, \texttt{.target},
  \texttt{.address\_size}), functions and kernels (\texttt{.func},
  \texttt{.entry}), and variables in a state space (\texttt{.reg},
  \texttt{.global}, \texttt{.shared}).  A \term{label}[ラベル] is a
  name followed by a colon in front of a statement; branch instructions
  use it as their target.
\end{definition}

\begin{example}[A complete PTX kernel]
  \label{ex:l10-ptx-kernel}
  The kernel below doubles one element of a global array per thread.
\begin{lstlisting}[language=PTX, numbers=none]
.version 8.0
.target sm_80
.address_size 64
.global .f32 buf[1024];

.visible .entry twice()
{
    .reg .b32  %r<2>;
    .reg .b64  %rd<4>;
    .reg .f32  %f<3>;
start:
    mov.u32        %r1, %tid.x;
    mov.u64        %rd1, buf;      // address of buf
    mul.wide.u32   %rd2, %r1, 4;   // byte offset 4*tid
    add.u64        %rd3, %rd1, %rd2;
    ld.global.f32  %f1, [%rd3];
    add.f32        %f2, %f1, %f1;  // 2 * buf[tid]
    st.global.f32  [%rd3], %f2;
    ret;
}
\end{lstlisting}
  The first four lines are directives; \texttt{.reg .b32 \%r<2>}
  declares the registers \texttt{\%r0} and \texttt{\%r1}.  The label
  \texttt{start} marks the first instruction, and every arithmetic and
  memory instruction names its type.
\end{example}

\margin{\texttt{\%r<2>} is PTX shorthand for a numbered range of
registers; \texttt{mul.wide.u32} multiplies two 32-bit values into a
64-bit result, as needed for a byte offset into a 64-bit address
space.}

\subsection{Other instructions}

Beyond arithmetic and memory access, PTX has the usual groups of
instructions.  Control flow uses \texttt{bra} (branch to a label),
\texttt{call} (call a function) and \texttt{ret} (return).\caution{There
is no stack-pointer special register in the PTX ISA, although such a
register is sometimes listed with \texttt{\%tid} and \texttt{\%ntid}.
Fixed-size local storage is a \texttt{.local} array, and
compiler-generated PTX addresses it through ordinary registers it names
\texttt{\%SP} and \texttt{\%SPL}.  Stack storage whose size is known
only at run time is allocated with \texttt{alloca}; \texttt{stacksave}
and \texttt{stackrestore} read and reset the stack pointer
\cite[§7.3]{nvidia2026ptx}.}
Synchronisation includes the barrier \texttt{bar.sync} that
\texttt{\_\_syncthreads()} compiles to and the
\term{memory fence}[メモリフェンス] instructions \texttt{membar} and
\texttt{fence}, which \enquote{enforce an ordering of memory
operations} \cite[§9.7.15.4]{nvidia2026ptx}: the memory accesses of a
thread before the fence become visible to other threads before its
accesses after the fence.  Atomic read-modify-write operations use
\texttt{atom}, always with an operation and a type, as in
\texttt{atom.global.add.u32}.\margin{Not every combination exists:
the half-precision atomic add is written \texttt{atom.add.noftz.f16}
(no flush to zero) and requires PTX ISA~6.3 and
\texttt{sm\_70} \cite[§9.7.15.5]{nvidia2026ptx}.}

Special registers (Lesson~4) give each thread its coordinates:
\texttt{\%tid} (thread index in the CTA), \texttt{\%ntid} (number of
threads per CTA dimension), \texttt{\%ctaid} and \texttt{\%nctaid}
(block index and grid size) \cite[§10]{nvidia2026ptx}.  The full list of instructions and
their variants is in the PTX ISA itself.

\section{IR and basic blocks}
\label{sec:l10-ir}

An optimizer needs the program in a form whose instructions are simple
and whose control flow is explicit.\source{M10, IR and basic block;
Aho et al.\ §6.2, §8.4 \cite{aho2006compilers}.}  The first requirement
leads to three-address code, the second to basic blocks and flow
graphs.

\subsection{Three-address code}

\begin{definition}[Three-address code]
  In \term{three-address code}[3番地コード] every instruction has at
  most one operator and at most three addresses (names of variables,
  temporaries or constants), in the form \texttt{A = B op C}.  A
  complex expression is broken into a sequence of such instructions with
  compiler-generated temporaries.
\end{definition}

LLVM IR and PTX are close to this form: most of their instructions
have one operation, one destination and one or two sources (PTX also
has three-source instructions such as \texttt{mad}).  In LLVM IR
\begin{lstlisting}[language=C++, numbers=none]
%sum = add i32 %x, %y
\end{lstlisting}
\texttt{\%sum} is the destination, \texttt{add} the operation,
\texttt{i32} the type of the result (a 32-bit integer) and
\texttt{\%x}, \texttt{\%y} the sources
\cite{llvm2026langref}.\margin{LLVM IR has one \texttt{add} for
signed and unsigned integers, since both produce the same bits; PTX
writes the interpretation into the opcode (\texttt{add.s32},
\texttt{add.u32}).}  The PTX form of the same operation is
\texttt{add.s32 \%r3, \%r1, \%r2}.  The expression
\texttt{y = a * b + c} becomes two instructions, \texttt{t = a * b}
and \texttt{y = t + c}.

\subsection{Basic blocks}

Lesson~6 introduced the \term*{basic block} informally.  The formal
definition is the unit on which most optimisations work.

\begin{definition}[Basic block]
  \label{def:l10-basic-block}
  A basic block\index{basic block} is a maximal sequence of consecutive
  instructions with one entry and one exit: control can enter only at
  its first instruction, and once it has entered, all instructions of
  the block execute in order, leaving the block only after the last
  one.  The last instruction is usually a control-flow instruction (a
  branch, a jump or a return); a block also ends where the next
  instruction is a branch target.
\end{definition}

\begin{definition}[Control flow graph, predecessor, successor]
  A \term{control flow graph (CFG)}[制御フローグラフ] has the basic
  blocks of a function as its nodes and an edge $B\to C$ whenever
  control can pass from the end of $B$ to the start of $C$: by a jump,
  by a taken conditional branch, or by falling through to the next
  block.\index{CFG|see{control flow graph (CFG)}}  $B$ is then a
  \term{predecessor}[先行ブロック] of $C$, and $C$ a
  \term{successor}[後続ブロック] of $B$.  An \emph{entry} node
  precedes the block of the first instruction, and an \emph{exit}
  node follows every block that leaves the function.
\end{definition}

\begin{marginfigure}
  \centering
  % The diamond control-flow graph of an if/else (Lesson 10).
% B1 holds the condition and branches to B2 (then) and B3 (else); both
% fall into the join block B4.
% Redraws the CFG of outline figure m10l3_flowgraph.
% Usage: % The diamond control-flow graph of an if/else (Lesson 10).
% B1 holds the condition and branches to B2 (then) and B3 (else); both
% fall into the join block B4.
% Redraws the CFG of outline figure m10l3_flowgraph.
% Usage: % The diamond control-flow graph of an if/else (Lesson 10).
% B1 holds the condition and branches to B2 (then) and B3 (else); both
% fall into the join block B4.
% Redraws the CFG of outline figure m10l3_flowgraph.
% Usage: \input{figures/l10-diamond}
% Width: margin (50 mm); the drawing is about 34 mm wide.
\begin{tikzpicture}[x=1mm, y=1mm, font=\scriptsize\sffamily,
    >={Stealth[length=1.6mm]},
    bb/.style={draw, fill=ln-box, minimum width=9mm, minimum height=5mm,
               inner sep=1.5pt}]
  \node[bb, fill=ln-accent!15] (b1) at (17,30) {B1};
  \node[bb] (b2) at (5,17) {B2};
  \node[bb] (b3) at (29,17) {B3};
  \node[bb] (b4) at (17,4) {B4};
  \draw[->] (b1) -- node[above left, inner sep=1pt]
    {\iflangja{真}{then}} (b2);
  \draw[->] (b1) -- node[above right, inner sep=1pt]
    {\iflangja{偽}{else}} (b3);
  \draw[->] (b2) -- (b4);
  \draw[->] (b3) -- (b4);
\end{tikzpicture}

% Width: margin (50 mm); the drawing is about 34 mm wide.
\begin{tikzpicture}[x=1mm, y=1mm, font=\scriptsize\sffamily,
    >={Stealth[length=1.6mm]},
    bb/.style={draw, fill=ln-box, minimum width=9mm, minimum height=5mm,
               inner sep=1.5pt}]
  \node[bb, fill=ln-accent!15] (b1) at (17,30) {B1};
  \node[bb] (b2) at (5,17) {B2};
  \node[bb] (b3) at (29,17) {B3};
  \node[bb] (b4) at (17,4) {B4};
  \draw[->] (b1) -- node[above left, inner sep=1pt]
    {\iflangja{真}{then}} (b2);
  \draw[->] (b1) -- node[above right, inner sep=1pt]
    {\iflangja{偽}{else}} (b3);
  \draw[->] (b2) -- (b4);
  \draw[->] (b3) -- (b4);
\end{tikzpicture}

% Width: margin (50 mm); the drawing is about 34 mm wide.
\begin{tikzpicture}[x=1mm, y=1mm, font=\scriptsize\sffamily,
    >={Stealth[length=1.6mm]},
    bb/.style={draw, fill=ln-box, minimum width=9mm, minimum height=5mm,
               inner sep=1.5pt}]
  \node[bb, fill=ln-accent!15] (b1) at (17,30) {B1};
  \node[bb] (b2) at (5,17) {B2};
  \node[bb] (b3) at (29,17) {B3};
  \node[bb] (b4) at (17,4) {B4};
  \draw[->] (b1) -- node[above left, inner sep=1pt]
    {\iflangja{真}{then}} (b2);
  \draw[->] (b1) -- node[above right, inner sep=1pt]
    {\iflangja{偽}{else}} (b3);
  \draw[->] (b2) -- (b4);
  \draw[->] (b3) -- (b4);
\end{tikzpicture}

  \caption{The diamond CFG of an \texttt{if}/\texttt{else}.}
  \label{fig:l10-diamond}
\end{marginfigure}
An \texttt{if}/\texttt{else} gives the diamond of Lesson~6
(\cref{fig:l10-diamond}): the block B1 with the condition has two
successors, the two branches B2 and B3, which share one successor, the
join block B4; the edges are B1$\to$B2, B1$\to$B3, B2$\to$B4 and
B3$\to$B4.  Optimisation inside one block is simple, because there is
only one path through it: if a block computes \texttt{a + b} twice and
neither operand is redefined in between, the second computation can
reuse the result of the first.

\begin{definition}[Local optimization]
  A \term{local optimization}[局所最適化] considers one basic block at a
  time, with no information about the paths that lead into the block or
  out of it.
\end{definition}

\subsection{Finding the basic blocks}

The blocks of a list of instructions follow from its leaders.

\begin{definition}[Leader]
  A \term{leader}[リーダ（先頭命令）] is an instruction that starts a
  basic block.  Three rules find all leaders: (1) the first instruction
  is a leader; (2) every target of a conditional or unconditional jump
  is a leader; (3) every instruction that immediately follows a jump is
  a leader.  Each basic block then runs from one leader up to, and not
  including, the next leader.
\end{definition}

\margin{The reconvergence point of a divergent branch (Lesson~6) is
always a leader: it is the first instruction of the join block, a
branch target.}

\begin{example}[Basic blocks of a PTX loop]
  \label{ex:l10-leaders}
  The loop below sums $0+1+\cdots+(n-1)$, with $i$ in \texttt{\%r1},
  the sum $s$ in \texttt{\%r2} and $n$ in \texttt{\%r3}.
\begin{lstlisting}[language=PTX, numbers=none]
      mov.u32  %r1, 0;             // (1) i = 0
      mov.u32  %r2, 0;             //     s = 0
loop: setp.hs.u32 %p1, %r1, %r3;   // (2) i >= n (unsigned)
@%p1  bra      done;
      add.u32  %r2, %r2, %r1;      // (3) s = s + i
      add.u32  %r1, %r1, 1;        //     i = i + 1
      bra      loop;
done: st.global.u32 [%rd1], %r2;   // (2),(3)
\end{lstlisting}
  The comments mark the rules that make an instruction a leader.  The
  first \texttt{mov} is a leader by rule~1.  The \texttt{setp} at
  \texttt{loop} is the target of \texttt{bra loop} (rule~2).  The first
  \texttt{add} follows the conditional \texttt{@\%p1 bra done}
  (rule~3).\caution{A guarded branch \texttt{@\%p1 bra} is a conditional
  jump: the instruction after it is a leader even though it is also
  reached by simply falling through.}  The store at \texttt{done} is both a target and the
  instruction after \texttt{bra loop}.  The four leaders give four
  blocks: B1 (the two \texttt{mov}s), B2 (\texttt{setp} and the
  conditional branch), B3 (the two \texttt{add}s and \texttt{bra
  loop}) and B4 (the store).  \Cref{fig:l10-loop-cfg} shows the CFG:
  B1$\to$B2; B2$\to$B3 by falling through and B2$\to$B4 when the
  branch is taken; and B3$\to$B2, the edge that closes the loop.
\end{example}

\begin{figure}[tb]
  \centering
  % Basic blocks and control-flow graph of a PTX counting loop (Lesson 10).
% B1 initialises i and s; B2 tests i >= n and branches to done; B3 adds
% and jumps back to loop (back edge B3 -> B2); B4 stores the sum.
% Original (the outline figure m10l3_egptx2 shows an if-else; this is the
% book's own loop example).
% Usage: % Basic blocks and control-flow graph of a PTX counting loop (Lesson 10).
% B1 initialises i and s; B2 tests i >= n and branches to done; B3 adds
% and jumps back to loop (back edge B3 -> B2); B4 stores the sum.
% Original (the outline figure m10l3_egptx2 shows an if-else; this is the
% book's own loop example).
% Usage: % Basic blocks and control-flow graph of a PTX counting loop (Lesson 10).
% B1 initialises i and s; B2 tests i >= n and branches to done; B3 adds
% and jumps back to loop (back edge B3 -> B2); B4 stores the sum.
% Original (the outline figure m10l3_egptx2 shows an if-else; this is the
% book's own loop example).
% Usage: \input{figures/l10-loop-cfg}
% Width: body (120 mm); the drawing is about 104 mm wide.
\begin{tikzpicture}[x=1mm, y=1mm, font=\scriptsize\sffamily,
    >={Stealth[length=1.6mm]},
    bb/.style={draw, fill=ln-box, align=left, inner sep=2pt,
               font=\scriptsize\ttfamily},
    nm/.style={font=\scriptsize\sffamily\bfseries, inner sep=1pt}]
  \node[text=ln-muted] (en) at (35,58) {ENTRY};
  \node[bb] (b1) at (35,50) {mov.u32 \%r1, 0;\\mov.u32 \%r2, 0;};
  \node[bb, fill=ln-accent!15] (b2) at (35,37)
    {loop: setp.hs.u32 \%p1, \%r1, \%r3;\\\phantom{loop: }@\%p1 bra done;};
  \node[bb] (b3) at (35,20) {add.u32 \%r2, \%r2, \%r1;\\add.u32 \%r1, \%r1, 1;\\bra loop;};
  \node[bb] (b4) at (92,20) {done:\\st.global.u32 [\%rd1], \%r2;};
  \node[text=ln-muted] (ex) at (92,8) {EXIT};
  \node[nm, anchor=east] at (b1.west) {B1\,};
  \node[nm, anchor=south east] at (b2.north west) {B2};
  \node[nm, anchor=west] at (b3.east) {\,B3};
  \node[nm, anchor=west] at (b4.east) {\,B4};
  \draw[->] (en) -- (b1);
  \draw[->] (b1) -- (b2);
  \draw[->] (b2) -- node[right, inner sep=2pt] {\iflangja{フォールスルー（$i<n$）}{fall-through ($i<n$)}} (b3);
  \draw[->] (b2.east) -| node[pos=0.25, above, inner sep=1.5pt] {\iflangja{分岐成立（$i\ge n$）}{taken ($i\ge n$)}} (b4.north);
  \draw[->, ln-caution, thick] (b3.west) -- ++(-6,0) |- node[pos=0.25, left, inner sep=1.5pt]
    {\iflangja{後退辺}{back edge}} (b2.west);
  \draw[->] (b4) -- (ex);
\end{tikzpicture}

% Width: body (120 mm); the drawing is about 104 mm wide.
\begin{tikzpicture}[x=1mm, y=1mm, font=\scriptsize\sffamily,
    >={Stealth[length=1.6mm]},
    bb/.style={draw, fill=ln-box, align=left, inner sep=2pt,
               font=\scriptsize\ttfamily},
    nm/.style={font=\scriptsize\sffamily\bfseries, inner sep=1pt}]
  \node[text=ln-muted] (en) at (35,58) {ENTRY};
  \node[bb] (b1) at (35,50) {mov.u32 \%r1, 0;\\mov.u32 \%r2, 0;};
  \node[bb, fill=ln-accent!15] (b2) at (35,37)
    {loop: setp.hs.u32 \%p1, \%r1, \%r3;\\\phantom{loop: }@\%p1 bra done;};
  \node[bb] (b3) at (35,20) {add.u32 \%r2, \%r2, \%r1;\\add.u32 \%r1, \%r1, 1;\\bra loop;};
  \node[bb] (b4) at (92,20) {done:\\st.global.u32 [\%rd1], \%r2;};
  \node[text=ln-muted] (ex) at (92,8) {EXIT};
  \node[nm, anchor=east] at (b1.west) {B1\,};
  \node[nm, anchor=south east] at (b2.north west) {B2};
  \node[nm, anchor=west] at (b3.east) {\,B3};
  \node[nm, anchor=west] at (b4.east) {\,B4};
  \draw[->] (en) -- (b1);
  \draw[->] (b1) -- (b2);
  \draw[->] (b2) -- node[right, inner sep=2pt] {\iflangja{フォールスルー（$i<n$）}{fall-through ($i<n$)}} (b3);
  \draw[->] (b2.east) -| node[pos=0.25, above, inner sep=1.5pt] {\iflangja{分岐成立（$i\ge n$）}{taken ($i\ge n$)}} (b4.north);
  \draw[->, ln-caution, thick] (b3.west) -- ++(-6,0) |- node[pos=0.25, left, inner sep=1.5pt]
    {\iflangja{後退辺}{back edge}} (b2.west);
  \draw[->] (b4) -- (ex);
\end{tikzpicture}

% Width: body (120 mm); the drawing is about 104 mm wide.
\begin{tikzpicture}[x=1mm, y=1mm, font=\scriptsize\sffamily,
    >={Stealth[length=1.6mm]},
    bb/.style={draw, fill=ln-box, align=left, inner sep=2pt,
               font=\scriptsize\ttfamily},
    nm/.style={font=\scriptsize\sffamily\bfseries, inner sep=1pt}]
  \node[text=ln-muted] (en) at (35,58) {ENTRY};
  \node[bb] (b1) at (35,50) {mov.u32 \%r1, 0;\\mov.u32 \%r2, 0;};
  \node[bb, fill=ln-accent!15] (b2) at (35,37)
    {loop: setp.hs.u32 \%p1, \%r1, \%r3;\\\phantom{loop: }@\%p1 bra done;};
  \node[bb] (b3) at (35,20) {add.u32 \%r2, \%r2, \%r1;\\add.u32 \%r1, \%r1, 1;\\bra loop;};
  \node[bb] (b4) at (92,20) {done:\\st.global.u32 [\%rd1], \%r2;};
  \node[text=ln-muted] (ex) at (92,8) {EXIT};
  \node[nm, anchor=east] at (b1.west) {B1\,};
  \node[nm, anchor=south east] at (b2.north west) {B2};
  \node[nm, anchor=west] at (b3.east) {\,B3};
  \node[nm, anchor=west] at (b4.east) {\,B4};
  \draw[->] (en) -- (b1);
  \draw[->] (b1) -- (b2);
  \draw[->] (b2) -- node[right, inner sep=2pt] {\iflangja{フォールスルー（$i<n$）}{fall-through ($i<n$)}} (b3);
  \draw[->] (b2.east) -| node[pos=0.25, above, inner sep=1.5pt] {\iflangja{分岐成立（$i\ge n$）}{taken ($i\ge n$)}} (b4.north);
  \draw[->, ln-caution, thick] (b3.west) -- ++(-6,0) |- node[pos=0.25, left, inner sep=1.5pt]
    {\iflangja{後退辺}{back edge}} (b2.west);
  \draw[->] (b4) -- (ex);
\end{tikzpicture}

  \caption{The basic blocks and the CFG of the loop in
    \cref{ex:l10-leaders}.  The edge B3$\to$B2 returns to an earlier
    block and closes the loop.}
  \label{fig:l10-loop-cfg}
\end{figure}

An edge such as B3$\to$B2, which leads back
to a block that starts the loop, is called a \emph{back
edge}.\margin{The edge from the end of a loop body to the loop test is
a back edge; it is what makes the data-flow iteration of
\cref{sec:l10-rd-example} need more than one pass.}

\section{Introduction to data-flow analysis}
\label{sec:l10-dfa}

Many optimisations need facts that span several basic
blocks.\source{M10, introduction to data-flow analysis; Aho et al.\
§9.1--9.2 \cite{aho2006compilers}.}

\begin{definition}[Global optimization]
  A \term{global optimization}[大域最適化] uses information collected
  across the basic blocks of a whole function, taking into account all
  paths of its CFG.
\end{definition}

Every optimisation, local or global, must preserve the semantics of the
original program: for every input the optimised program produces the
same observable results.\caution{Floating point bends this rule by
default: nvcc contracts a multiply followed by an add into one fused
multiply-add, which rounds once instead of twice, so results can differ
in the last bit; \texttt{-{}-fmad=false} turns the contraction off
(NVCC §4.2.7.11 \cite{nvidia2026nvcc}).}  The common global optimisations, all treated
in Lesson~11, are the removal of redundant instructions (an expression
computed twice, on every path, with the same operands), constant and
\term{copy propagation}[コピー伝播] (replacing a variable by the
constant or the other variable it was assigned from, which often
exposes further redundancy), dead code elimination (removing
instructions whose results are never used), code motion (moving an
instruction to where it executes less often, e.g.\ out of a loop), and
the detection of induction variables with strength
reduction.\margin{Copy propagation in one line: after \texttt{y = x},
a later \texttt{z = y + 1} becomes \texttt{z = x + 1}; if no other use
of \texttt{y} remains, \texttt{y = x} is dead.}  None of them can be
decided inside one block: whether \texttt{x} in some block is still
the constant assigned three blocks earlier depends on every path in
between.

\begin{definition}[Data-flow analysis]
  The family of techniques called
  \term{data-flow analysis}[データフロー解析] computes, for every
  point of a program, a fact about all executions that can reach that
  point, by analysing the effect of each basic block and propagating
  the results along the edges of the CFG.  The analyses differ in the
  property they examine.
\end{definition}

\subsection{The abstraction: points, states and transfer functions}

Data-flow analysis describes a program by the states between its
statements.

\begin{definition}[Program point, transfer function]
  A \term{program point}[プログラム点] lies before or after a
  statement.  The analysis attaches to each point a state, the set of
  facts that hold there.  For a block $B$, $\mathrm{IN}[B]$ is the
  state at the point before its first statement and $\mathrm{OUT}[B]$
  the state after its last.  The \term{transfer function}[伝達関数]
  $f_s$ of a statement $s$ maps the state before $s$ to the state after
  it.
\end{definition}

Inside a block, the point after one statement is the point before the
next (\cref{fig:l10-transfer}), so the transfer function of a block of
statements $s_1,\dots,s_n$ is the composition
\begin{equation}
  \label{eq:l10-compose}
  \mathrm{OUT}[B] = f_B(\mathrm{IN}[B]),\qquad
  f_B = f_{s_n}\circ\cdots\circ f_{s_2}\circ f_{s_1}.
\end{equation}
Between blocks, the state on entry to a block must combine the states
of all its predecessors, since control can arrive from any of them.

\begin{figure}[tb]
  \centering
  % Program points and transfer functions (Lesson 10).  Left: a block B of
% three statements; a program point lies before and after each statement,
% IN[B] is the first and OUT[B] the last, and f_B is the composition of
% the statement transfer functions.  Right: at a join, IN[B] is the
% union (meet) of OUT of all predecessors (forward problem).
% Redraws the outline figure m10l4_trans with a generic block.
% Usage: % Program points and transfer functions (Lesson 10).  Left: a block B of
% three statements; a program point lies before and after each statement,
% IN[B] is the first and OUT[B] the last, and f_B is the composition of
% the statement transfer functions.  Right: at a join, IN[B] is the
% union (meet) of OUT of all predecessors (forward problem).
% Redraws the outline figure m10l4_trans with a generic block.
% Usage: % Program points and transfer functions (Lesson 10).  Left: a block B of
% three statements; a program point lies before and after each statement,
% IN[B] is the first and OUT[B] the last, and f_B is the composition of
% the statement transfer functions.  Right: at a join, IN[B] is the
% union (meet) of OUT of all predecessors (forward problem).
% Redraws the outline figure m10l4_trans with a generic block.
% Usage: \input{figures/l10-transfer}
% Width: body (120 mm); the drawing is about 112 mm wide.
\begin{tikzpicture}[x=1mm, y=1mm, font=\scriptsize\sffamily,
    >={Stealth[length=1.5mm]},
    st/.style={draw, fill=ln-box, minimum width=22mm, minimum height=5.5mm,
               inner sep=1pt},
    pt/.style={circle, fill=ln-source, inner sep=0pt, minimum size=1.6mm}]
  % left: one block
  \node[font=\scriptsize\sffamily\bfseries] at (20,47) {\iflangja{ブロック B の内部}{inside block B}};
  \draw[ln-muted, rounded corners=2pt] (5,3) rectangle (35,43);
  \foreach \k/\y in {1/34, 2/23, 3/12} {
    \node[st] (s\k) at (20,\y) {$s_\k$};
  }
  \foreach \y/\lab in {40.5/{$\mathrm{IN}[B]$}, 28.5/{}, 17.5/{}, 6.5/{$\mathrm{OUT}[B]$}} {
    \node[pt] at (20,\y) {};
  }
  \node[anchor=west, inner sep=1pt] at (37,40.5) {$\mathrm{IN}[B]$ \iflangja{（$s_1$ の前）}{(before $s_1$)}};
  \node[anchor=west, inner sep=1pt] at (37,28.5) {\iflangja{$s_1$ の後 $=$ $s_2$ の前}{after $s_1$ = before $s_2$}};
  \node[anchor=west, inner sep=1pt] at (37,17.5) {\iflangja{$s_2$ の後 $=$ $s_3$ の前}{after $s_2$ = before $s_3$}};
  \node[anchor=west, inner sep=1pt] at (37,6.5) {$\mathrm{OUT}[B]$ \iflangja{（$s_3$ の後）}{(after $s_3$)}};
  \foreach \k/\y in {1/34, 2/23, 3/12}
    \node[anchor=east, text=ln-accent, inner sep=1pt] at (8.5,\y) {$f_{s_\k}$};
  \node[text=ln-muted, align=left, anchor=west] at (0,-2.5)
    {$f_B = f_{s_3}\circ f_{s_2}\circ f_{s_1}$};
  % right: a join
  \node[font=\scriptsize\sffamily\bfseries] at (94,47) {\iflangja{合流点}{at a join}};
  \node[st, minimum width=14mm] (p1) at (82,34) {$P_1$};
  \node[st, minimum width=14mm] (p2) at (106,34) {$P_2$};
  \node[st, minimum width=14mm] (bb) at (94,12) {$B$};
  \draw[->] (p1) -- node[anchor=east, xshift=-2.5mm, inner sep=0pt, pos=0.5] {$\mathrm{OUT}[P_1]$} (bb);
  \draw[->] (p2) -- node[anchor=west, xshift=2.5mm, inner sep=0pt, pos=0.5] {$\mathrm{OUT}[P_2]$} (bb);
  \node[anchor=north, inner sep=1pt] at (94,6)
    {$\mathrm{IN}[B]=\mathrm{OUT}[P_1]\cup\mathrm{OUT}[P_2]$};
\end{tikzpicture}

% Width: body (120 mm); the drawing is about 112 mm wide.
\begin{tikzpicture}[x=1mm, y=1mm, font=\scriptsize\sffamily,
    >={Stealth[length=1.5mm]},
    st/.style={draw, fill=ln-box, minimum width=22mm, minimum height=5.5mm,
               inner sep=1pt},
    pt/.style={circle, fill=ln-source, inner sep=0pt, minimum size=1.6mm}]
  % left: one block
  \node[font=\scriptsize\sffamily\bfseries] at (20,47) {\iflangja{ブロック B の内部}{inside block B}};
  \draw[ln-muted, rounded corners=2pt] (5,3) rectangle (35,43);
  \foreach \k/\y in {1/34, 2/23, 3/12} {
    \node[st] (s\k) at (20,\y) {$s_\k$};
  }
  \foreach \y/\lab in {40.5/{$\mathrm{IN}[B]$}, 28.5/{}, 17.5/{}, 6.5/{$\mathrm{OUT}[B]$}} {
    \node[pt] at (20,\y) {};
  }
  \node[anchor=west, inner sep=1pt] at (37,40.5) {$\mathrm{IN}[B]$ \iflangja{（$s_1$ の前）}{(before $s_1$)}};
  \node[anchor=west, inner sep=1pt] at (37,28.5) {\iflangja{$s_1$ の後 $=$ $s_2$ の前}{after $s_1$ = before $s_2$}};
  \node[anchor=west, inner sep=1pt] at (37,17.5) {\iflangja{$s_2$ の後 $=$ $s_3$ の前}{after $s_2$ = before $s_3$}};
  \node[anchor=west, inner sep=1pt] at (37,6.5) {$\mathrm{OUT}[B]$ \iflangja{（$s_3$ の後）}{(after $s_3$)}};
  \foreach \k/\y in {1/34, 2/23, 3/12}
    \node[anchor=east, text=ln-accent, inner sep=1pt] at (8.5,\y) {$f_{s_\k}$};
  \node[text=ln-muted, align=left, anchor=west] at (0,-2.5)
    {$f_B = f_{s_3}\circ f_{s_2}\circ f_{s_1}$};
  % right: a join
  \node[font=\scriptsize\sffamily\bfseries] at (94,47) {\iflangja{合流点}{at a join}};
  \node[st, minimum width=14mm] (p1) at (82,34) {$P_1$};
  \node[st, minimum width=14mm] (p2) at (106,34) {$P_2$};
  \node[st, minimum width=14mm] (bb) at (94,12) {$B$};
  \draw[->] (p1) -- node[anchor=east, xshift=-2.5mm, inner sep=0pt, pos=0.5] {$\mathrm{OUT}[P_1]$} (bb);
  \draw[->] (p2) -- node[anchor=west, xshift=2.5mm, inner sep=0pt, pos=0.5] {$\mathrm{OUT}[P_2]$} (bb);
  \node[anchor=north, inner sep=1pt] at (94,6)
    {$\mathrm{IN}[B]=\mathrm{OUT}[P_1]\cup\mathrm{OUT}[P_2]$};
\end{tikzpicture}

% Width: body (120 mm); the drawing is about 112 mm wide.
\begin{tikzpicture}[x=1mm, y=1mm, font=\scriptsize\sffamily,
    >={Stealth[length=1.5mm]},
    st/.style={draw, fill=ln-box, minimum width=22mm, minimum height=5.5mm,
               inner sep=1pt},
    pt/.style={circle, fill=ln-source, inner sep=0pt, minimum size=1.6mm}]
  % left: one block
  \node[font=\scriptsize\sffamily\bfseries] at (20,47) {\iflangja{ブロック B の内部}{inside block B}};
  \draw[ln-muted, rounded corners=2pt] (5,3) rectangle (35,43);
  \foreach \k/\y in {1/34, 2/23, 3/12} {
    \node[st] (s\k) at (20,\y) {$s_\k$};
  }
  \foreach \y/\lab in {40.5/{$\mathrm{IN}[B]$}, 28.5/{}, 17.5/{}, 6.5/{$\mathrm{OUT}[B]$}} {
    \node[pt] at (20,\y) {};
  }
  \node[anchor=west, inner sep=1pt] at (37,40.5) {$\mathrm{IN}[B]$ \iflangja{（$s_1$ の前）}{(before $s_1$)}};
  \node[anchor=west, inner sep=1pt] at (37,28.5) {\iflangja{$s_1$ の後 $=$ $s_2$ の前}{after $s_1$ = before $s_2$}};
  \node[anchor=west, inner sep=1pt] at (37,17.5) {\iflangja{$s_2$ の後 $=$ $s_3$ の前}{after $s_2$ = before $s_3$}};
  \node[anchor=west, inner sep=1pt] at (37,6.5) {$\mathrm{OUT}[B]$ \iflangja{（$s_3$ の後）}{(after $s_3$)}};
  \foreach \k/\y in {1/34, 2/23, 3/12}
    \node[anchor=east, text=ln-accent, inner sep=1pt] at (8.5,\y) {$f_{s_\k}$};
  \node[text=ln-muted, align=left, anchor=west] at (0,-2.5)
    {$f_B = f_{s_3}\circ f_{s_2}\circ f_{s_1}$};
  % right: a join
  \node[font=\scriptsize\sffamily\bfseries] at (94,47) {\iflangja{合流点}{at a join}};
  \node[st, minimum width=14mm] (p1) at (82,34) {$P_1$};
  \node[st, minimum width=14mm] (p2) at (106,34) {$P_2$};
  \node[st, minimum width=14mm] (bb) at (94,12) {$B$};
  \draw[->] (p1) -- node[anchor=east, xshift=-2.5mm, inner sep=0pt, pos=0.5] {$\mathrm{OUT}[P_1]$} (bb);
  \draw[->] (p2) -- node[anchor=west, xshift=2.5mm, inner sep=0pt, pos=0.5] {$\mathrm{OUT}[P_2]$} (bb);
  \node[anchor=north, inner sep=1pt] at (94,6)
    {$\mathrm{IN}[B]=\mathrm{OUT}[P_1]\cup\mathrm{OUT}[P_2]$};
\end{tikzpicture}

  \caption{Left: the program points of a block of three statements;
    each statement's transfer function maps the state before it to the
    state after it.  Right: at a join, the meet (here the union)
    combines the states of the predecessors.}
  \label{fig:l10-transfer}
\end{figure}

\begin{definition}[Meet operator, forward analysis]
  The \term{meet operator}[合流演算子] combines the states that arrive
  at a join from several edges into one state.  In a
  \term{forward analysis}[前向き解析] the states flow in the direction
  of control: the equations are
  \begin{equation}
    \label{eq:l10-forward}
    \mathrm{IN}[B] = \bigcup_{P\in\mathrm{pred}(B)}\mathrm{OUT}[P],\qquad
    \mathrm{OUT}[B] = f_B(\mathrm{IN}[B]).
  \end{equation}
\end{definition}

A backward problem runs the other way, from the uses of a value back
to where it was computed: $\mathrm{IN}[B]=f_B(\mathrm{OUT}[B])$ and
$\mathrm{OUT}[B]=\bigcup_{S\in\mathrm{succ}(B)}\mathrm{IN}[S]$.\margin{Liveness,
the backward analysis behind register allocation, is the subject of
Lesson~11.  Both analyses of this course use union as the meet.}  In
the CFG of \cref{fig:l10-rd-cfg} below, B5 has the predecessors B3 and
B4, so in a forward analysis
$\mathrm{IN}[\mathrm{B5}]=\mathrm{OUT}[\mathrm{B3}]\cup\mathrm{OUT}[\mathrm{B4}]$;
B5 has the successors B2 and EXIT, so in a backward analysis
$\mathrm{OUT}[\mathrm{B5}]=\mathrm{IN}[\mathrm{B2}]\cup\mathrm{IN}[\mathrm{EXIT}]$.

\subsection{Reaching definitions}

The first analysis concerns where values come from.

\begin{definition}[Definition, use]
  A \term{definition}[定義] of a variable is a statement that assigns
  it a value; a \term{use}[使用] is a statement that reads it.  In
  \texttt{a = x + y} the statement defines \texttt{a} and uses
  \texttt{x} and \texttt{y}.
\end{definition}

\begin{definition}[Reaching definitions]
  A definition $d$ \emph{reaches} a point $p$ if there is a path from
  the point immediately after $d$ to $p$ along which $d$ is not killed,
  that is, along which no other definition of the same variable
  overwrites it.  The analysis of \term{reaching definitions}[到達定義]
  computes, for every program point, the set of definitions that reach
  it.
\end{definition}

\caution{One path suffices: a definition reaches $p$ if it survives on
\emph{some} path to $p$, not necessarily on all of them.  The analysis
is conservative in this direction, since it may list a definition
whose path is never taken at run time.}  If exactly one definition of
\texttt{x}, say \texttt{x = 4}, reaches a use of \texttt{x}, the use
can be replaced by the constant~4; if a second definition also reaches
it, it cannot.  This is how constant propagation uses the analysis.

\begin{definition}[Gen and kill sets]
  For a definition $d$\,: \texttt{u = v + w}, the
  \term{gen set}[gen集合] is $\mathrm{gen}_d=\{d\}$, the definition the
  statement creates, and the \term{kill set}[kill集合]
  $\mathrm{kill}_d$ is the set of all other definitions of \texttt{u}
  in the program, which the statement overwrites.  The transfer
  function of $d$ is
  \begin{equation}
    \label{eq:l10-fd}
    f_d(x) = \mathrm{gen}_d \cup (x - \mathrm{kill}_d).
  \end{equation}
\end{definition}

\margin{Notation: $x - y$ is set difference, the elements of $x$ not in
$y$; $x$ stands for the set of definitions that reach the point before
the statement.}  A definition that arrives at $d$ survives unless $d$
kills it, and $d$ itself is added.  For a block, the functions compose
as in \cref{eq:l10-compose}.  Two statements give
\begin{align*}
  f_2(f_1(x)) &= \mathrm{gen}_2 \cup \bigl(\bigl(\mathrm{gen}_1 \cup
    (x - \mathrm{kill}_1)\bigr) - \mathrm{kill}_2\bigr)\\
  &= \bigl(\mathrm{gen}_2 \cup (\mathrm{gen}_1 - \mathrm{kill}_2)\bigr)
     \cup \bigl(x - (\mathrm{kill}_1 \cup \mathrm{kill}_2)\bigr),
\end{align*}
which has the same form as \cref{eq:l10-fd}.  By induction, a block of
$n$ statements has the transfer function
$f_B(x)=\mathrm{gen}_B\cup(x-\mathrm{kill}_B)$ with
\begin{equation}
  \label{eq:l10-genkill}
  \begin{aligned}
    \mathrm{kill}_B &= \mathrm{kill}_1\cup\mathrm{kill}_2\cup\cdots\cup\mathrm{kill}_n,\\
    \mathrm{gen}_B &= \mathrm{gen}_n \cup (\mathrm{gen}_{n-1}-\mathrm{kill}_n)
      \cup (\mathrm{gen}_{n-2}-\mathrm{kill}_{n-1}-\mathrm{kill}_n)\\
      &\qquad\cup\cdots\cup(\mathrm{gen}_1-\mathrm{kill}_2-\cdots-\mathrm{kill}_n).
  \end{aligned}
\end{equation}
A block kills every definition that one of its statements kills, and
generates every definition of its own that no later statement in the
block overwrites.\tip{Read a block bottom-up to find
$\mathrm{gen}_B$: a definition belongs to it unless a later statement
of the block defines the same variable again.}

\begin{example}[Gen and kill sets of a CFG]
  \label{ex:l10-genkill}
  The CFG of \cref{fig:l10-rd-cfg} sums the absolute values of
  \texttt{x[0]} to \texttt{x[n-1]} into \texttt{s} (for $n\ge 1$;
  the test is at the bottom of the loop).  It has seven
  definitions: \texttt{i} is defined by $d_1$ and $d_7$, \texttt{s} by
  $d_2$, $d_4$ and $d_6$, \texttt{t} by $d_3$ and $d_5$.  Block B4
  holds two statements: $d_5$\,: \texttt{t = -t} with
  $\mathrm{gen}=\{d_5\}$, $\mathrm{kill}=\{d_3\}$, and
  $d_6$\,: \texttt{s = s + t} with $\mathrm{gen}=\{d_6\}$,
  $\mathrm{kill}=\{d_2,d_4\}$.  By \cref{eq:l10-genkill},
  $\mathrm{kill}_{\mathrm{B4}}=\{d_2,d_3,d_4\}$ and
  $\mathrm{gen}_{\mathrm{B4}}=\{d_6\}\cup(\{d_5\}-\{d_2,d_4\})=\{d_5,d_6\}$.
  \Cref{tab:l10-genkill} lists all five blocks.
\end{example}

\begin{figure}[tb]
  \centering
  % The control-flow graph of the reaching-definitions example (Lesson 10).
% A loop over x[0..n-1] that sums absolute values: B1 initialises i and s,
% B2 loads t, B3 adds a positive t, B4 negates t and adds it, B5
% increments i and loops back to B2 or leaves to EXIT.  Definitions are
% numbered d1..d7.
% Original (the book's own CFG; not the outline's m10l4_gnk2 example).
% Usage: % The control-flow graph of the reaching-definitions example (Lesson 10).
% A loop over x[0..n-1] that sums absolute values: B1 initialises i and s,
% B2 loads t, B3 adds a positive t, B4 negates t and adds it, B5
% increments i and loops back to B2 or leaves to EXIT.  Definitions are
% numbered d1..d7.
% Original (the book's own CFG; not the outline's m10l4_gnk2 example).
% Usage: % The control-flow graph of the reaching-definitions example (Lesson 10).
% A loop over x[0..n-1] that sums absolute values: B1 initialises i and s,
% B2 loads t, B3 adds a positive t, B4 negates t and adds it, B5
% increments i and loops back to B2 or leaves to EXIT.  Definitions are
% numbered d1..d7.
% Original (the book's own CFG; not the outline's m10l4_gnk2 example).
% Usage: \input{figures/l10-rd-cfg}
% Width: body (120 mm); the drawing is about 92 mm wide.
\begin{tikzpicture}[x=1mm, y=1mm, font=\scriptsize\sffamily,
    >={Stealth[length=1.6mm]},
    bb/.style={draw, fill=ln-box, align=left, inner sep=2pt, font=\scriptsize},
    nm/.style={font=\scriptsize\sffamily\bfseries, inner sep=1pt}]
  \node[text=ln-muted] (en) at (45,66) {ENTRY};
  \node[bb] (b1) at (45,57) {$d_1$: \texttt{i = 0}\\$d_2$: \texttt{s = 0}};
  \node[bb, fill=ln-accent!15] (b2) at (45,42) {$d_3$: \texttt{t = x[i]}};
  \node[bb] (b3) at (22,27) {$d_4$: \texttt{s = s + t}};
  \node[bb] (b4) at (68,27) {$d_5$: \texttt{t = -t}\\$d_6$: \texttt{s = s + t}};
  \node[bb] (b5) at (45,12) {$d_7$: \texttt{i = i + 1}};
  \node[text=ln-muted] (ex) at (45,1) {EXIT};
  \node[nm, anchor=east] at (b1.west) {B1\,};
  \node[nm, anchor=east] at (b2.west) {B2\,};
  \node[nm, anchor=east] at (b3.west) {B3\,};
  \node[nm, anchor=west] at (b4.east) {\,B4};
  \node[nm, anchor=east] at (b5.west) {B5\,};
  \draw[->] (en) -- (b1);
  \draw[->] (b1) -- (b2);
  \draw[->] (b2) -- node[above left, inner sep=1pt] {\texttt{t > 0}} (b3);
  \draw[->] (b2) -- node[above right, inner sep=1pt] {\texttt{t <= 0}} (b4);
  \draw[->] (b3) -- (b5);
  \draw[->] (b4) -- (b5);
  \draw[->] (b5) -- node[right, inner sep=1.5pt] {\texttt{i >= n}} (ex);
  \draw[->, ln-caution, thick] (b5.east) -- ++(40,0) |- node[pos=0.25, right, inner sep=1.5pt]
    {\texttt{i < n}} (b2.east);
\end{tikzpicture}

% Width: body (120 mm); the drawing is about 92 mm wide.
\begin{tikzpicture}[x=1mm, y=1mm, font=\scriptsize\sffamily,
    >={Stealth[length=1.6mm]},
    bb/.style={draw, fill=ln-box, align=left, inner sep=2pt, font=\scriptsize},
    nm/.style={font=\scriptsize\sffamily\bfseries, inner sep=1pt}]
  \node[text=ln-muted] (en) at (45,66) {ENTRY};
  \node[bb] (b1) at (45,57) {$d_1$: \texttt{i = 0}\\$d_2$: \texttt{s = 0}};
  \node[bb, fill=ln-accent!15] (b2) at (45,42) {$d_3$: \texttt{t = x[i]}};
  \node[bb] (b3) at (22,27) {$d_4$: \texttt{s = s + t}};
  \node[bb] (b4) at (68,27) {$d_5$: \texttt{t = -t}\\$d_6$: \texttt{s = s + t}};
  \node[bb] (b5) at (45,12) {$d_7$: \texttt{i = i + 1}};
  \node[text=ln-muted] (ex) at (45,1) {EXIT};
  \node[nm, anchor=east] at (b1.west) {B1\,};
  \node[nm, anchor=east] at (b2.west) {B2\,};
  \node[nm, anchor=east] at (b3.west) {B3\,};
  \node[nm, anchor=west] at (b4.east) {\,B4};
  \node[nm, anchor=east] at (b5.west) {B5\,};
  \draw[->] (en) -- (b1);
  \draw[->] (b1) -- (b2);
  \draw[->] (b2) -- node[above left, inner sep=1pt] {\texttt{t > 0}} (b3);
  \draw[->] (b2) -- node[above right, inner sep=1pt] {\texttt{t <= 0}} (b4);
  \draw[->] (b3) -- (b5);
  \draw[->] (b4) -- (b5);
  \draw[->] (b5) -- node[right, inner sep=1.5pt] {\texttt{i >= n}} (ex);
  \draw[->, ln-caution, thick] (b5.east) -- ++(40,0) |- node[pos=0.25, right, inner sep=1.5pt]
    {\texttt{i < n}} (b2.east);
\end{tikzpicture}

% Width: body (120 mm); the drawing is about 92 mm wide.
\begin{tikzpicture}[x=1mm, y=1mm, font=\scriptsize\sffamily,
    >={Stealth[length=1.6mm]},
    bb/.style={draw, fill=ln-box, align=left, inner sep=2pt, font=\scriptsize},
    nm/.style={font=\scriptsize\sffamily\bfseries, inner sep=1pt}]
  \node[text=ln-muted] (en) at (45,66) {ENTRY};
  \node[bb] (b1) at (45,57) {$d_1$: \texttt{i = 0}\\$d_2$: \texttt{s = 0}};
  \node[bb, fill=ln-accent!15] (b2) at (45,42) {$d_3$: \texttt{t = x[i]}};
  \node[bb] (b3) at (22,27) {$d_4$: \texttt{s = s + t}};
  \node[bb] (b4) at (68,27) {$d_5$: \texttt{t = -t}\\$d_6$: \texttt{s = s + t}};
  \node[bb] (b5) at (45,12) {$d_7$: \texttt{i = i + 1}};
  \node[text=ln-muted] (ex) at (45,1) {EXIT};
  \node[nm, anchor=east] at (b1.west) {B1\,};
  \node[nm, anchor=east] at (b2.west) {B2\,};
  \node[nm, anchor=east] at (b3.west) {B3\,};
  \node[nm, anchor=west] at (b4.east) {\,B4};
  \node[nm, anchor=east] at (b5.west) {B5\,};
  \draw[->] (en) -- (b1);
  \draw[->] (b1) -- (b2);
  \draw[->] (b2) -- node[above left, inner sep=1pt] {\texttt{t > 0}} (b3);
  \draw[->] (b2) -- node[above right, inner sep=1pt] {\texttt{t <= 0}} (b4);
  \draw[->] (b3) -- (b5);
  \draw[->] (b4) -- (b5);
  \draw[->] (b5) -- node[right, inner sep=1.5pt] {\texttt{i >= n}} (ex);
  \draw[->, ln-caution, thick] (b5.east) -- ++(40,0) |- node[pos=0.25, right, inner sep=1.5pt]
    {\texttt{i < n}} (b2.east);
\end{tikzpicture}

  \caption{The CFG of the reaching-definitions example: a loop that
    adds \texttt{t} when it is positive and $-$\texttt{t} otherwise.
    The edge B5$\to$B2 is the back edge.}
  \label{fig:l10-rd-cfg}
\end{figure}

\begin{table}[tb]
  \centering
  \small
  \caption{Gen and kill sets of the blocks of \cref{fig:l10-rd-cfg}.}
  \label{tab:l10-genkill}
  \begin{tabular}{@{}l>{\raggedright\arraybackslash}p{24mm}
                  >{\raggedright\arraybackslash}p{24mm}@{}}
    \toprule
    Block & $\mathrm{gen}_B$ & $\mathrm{kill}_B$ \\
    \midrule
    B1 & $\{d_1,d_2\}$ & $\{d_4,d_6,d_7\}$ \\
    B2 & $\{d_3\}$ & $\{d_5\}$ \\
    B3 & $\{d_4\}$ & $\{d_2,d_6\}$ \\
    B4 & $\{d_5,d_6\}$ & $\{d_2,d_3,d_4\}$ \\
    B5 & $\{d_7\}$ & $\{d_1\}$ \\
    \bottomrule
  \end{tabular}
\end{table}

\section{An example of reaching definitions}
\label{sec:l10-rd-example}

With the gen and kill sets of every block, the reaching definitions
follow from a system of equations over the whole CFG.\source{M10,
example of reaching definitions; Aho et al.\ §9.2.4
\cite{aho2006compilers}.}

\subsection{Equations and the iterative algorithm}

For every block $B$ other than ENTRY, reaching definitions is the
forward problem of \cref{eq:l10-forward} with union as the meet:
\begin{equation}
  \label{eq:l10-rd}
  \mathrm{IN}[B]=\bigcup_{P\in\mathrm{pred}(B)}\mathrm{OUT}[P],\qquad
  \mathrm{OUT}[B]=\mathrm{gen}_B\cup(\mathrm{IN}[B]-\mathrm{kill}_B),
\end{equation}
together with the boundary condition
$\mathrm{OUT}[\mathrm{ENTRY}]=\emptyset$: no definition reaches the
start of the function.\margin{The exit node has no statements, so
$\mathrm{OUT}[\mathrm{EXIT}]=\mathrm{IN}[\mathrm{EXIT}]$, the
definitions that reach the end of the function.}  With a loop in the
CFG the equations are circular: $\mathrm{IN}[\mathrm{B2}]$ depends on
$\mathrm{OUT}[\mathrm{B5}]$, which depends on
$\mathrm{IN}[\mathrm{B2}]$.  They are solved by iteration.

\begin{definition}[Iterative algorithm, fixed point]
  The \term{iterative algorithm}[反復アルゴリズム] for a data-flow
  problem starts from an initial value for every state and evaluates
  the equations for all blocks in turn, pass after pass, until a whole
  pass changes no state.  The states are then a solution of the
  equations, a \term{fixed point}[不動点] of the passes: applying the
  equations once more leaves them unchanged.
\end{definition}

For reaching definitions the algorithm reads:
\begin{lstlisting}[mathescape, numbers=none]
$\lxOUT[\lxENTRY] = \emptyset$
for each block $B \ne \lxENTRY$:  $\lxOUT[B] = \emptyset$
while some $\lxOUT$ changed in the last pass:
    for each block $B \ne \lxENTRY$:
        $\lxIN[B] = \bigcup_{P \in \lxpred(B)} \lxOUT[P]$
        $\lxOUT[B] = \lxgen_B \cup (\lxIN[B] - \lxkill_B)$
\end{lstlisting}

The algorithm terminates.  Every $\mathrm{OUT}$ set starts empty, and
since \cref{eq:l10-rd} only adds definitions when its inputs grow, no
set ever shrinks; the sets are subsets of the finite set of
definitions, so they can grow only finitely often.\caution{The order in
which a pass visits the blocks changes the number of passes, not the
result.  Visiting predecessors before successors, as in the example
below, lets the information travel far in one pass; only back edges
then require another.}  Starting from empty sets yields the smallest
solution, which contains only definitions that really have a path to
the point.

\begin{definition}[Bit vector]
  A \term{bit vector}[ビットベクタ] represents a subset of a fixed,
  numbered set of $N$ elements as $N$ bits, bit $k$ being 1 when
  element $k$ is in the subset.  Union becomes bitwise OR, intersection
  bitwise AND and $x-y$ becomes $x$ AND NOT $y$.
\end{definition}

With the definitions numbered $d_1,\dots,d_N$, each set of the
algorithm is one bit vector, and a whole transfer function is two
bitwise operations per machine word.\margin{The seven definitions of the
example fit into one machine word; a function with 1000 definitions
needs 16 words of 64 bits per set.}

\subsection{The example worked out}

\Cref{tab:l10-rd-trace} applies the algorithm to the CFG of
\cref{fig:l10-rd-cfg}, visiting B1 to B5 and EXIT in this order.  Each
bit vector lists $d_1,\dots,d_7$ grouped by the block that contains the
definition (B1: $d_1d_2$; B2: $d_3$; B3: $d_4$; B4: $d_5d_6$; B5:
$d_7$).

\begin{table}[tb]
  \centering
  \small
  \setlength{\tabcolsep}{4pt}
  \caption{Reaching definitions for \cref{fig:l10-rd-cfg}, pass by
    pass.  Bit $k$ stands for $d_k$; the superscript is the pass.  A
    third pass reproduces pass~2.}
  \label{tab:l10-rd-trace}
  \begin{tabular}{@{}l*{5}{c}@{}}
    \toprule
    Block & $\mathrm{OUT}^0$ & $\mathrm{IN}^1$ & $\mathrm{OUT}^1$
      & $\mathrm{IN}^2$ & $\mathrm{OUT}^2$ \\
    \midrule
    B1 & \texttt{00 0 0 00 0} & \texttt{00 0 0 00 0} & \texttt{11 0 0 00 0}
       & \texttt{00 0 0 00 0} & \texttt{11 0 0 00 0} \\
    B2 & \texttt{00 0 0 00 0} & \texttt{11 0 0 00 0} & \texttt{11 1 0 00 0}
       & \texttt{11 1 1 11 1} & \texttt{11 1 1 01 1} \\
    B3 & \texttt{00 0 0 00 0} & \texttt{11 1 0 00 0} & \texttt{10 1 1 00 0}
       & \texttt{11 1 1 01 1} & \texttt{10 1 1 00 1} \\
    B4 & \texttt{00 0 0 00 0} & \texttt{11 1 0 00 0} & \texttt{10 0 0 11 0}
       & \texttt{11 1 1 01 1} & \texttt{10 0 0 11 1} \\
    B5 & \texttt{00 0 0 00 0} & \texttt{10 1 1 11 0} & \texttt{00 1 1 11 1}
       & \texttt{10 1 1 11 1} & \texttt{00 1 1 11 1} \\
    EXIT & \texttt{00 0 0 00 0} & \texttt{00 1 1 11 1} & \texttt{00 1 1 11 1}
       & \texttt{00 1 1 11 1} & \texttt{00 1 1 11 1} \\
    \bottomrule
  \end{tabular}
\end{table}

\begin{example}[The first pass]
  \label{ex:l10-pass1}
  B1 has only ENTRY as predecessor, so $\mathrm{IN}[\mathrm{B1}]=\emptyset$
  and $\mathrm{OUT}[\mathrm{B1}]=\mathrm{gen}_{\mathrm{B1}}=\{d_1,d_2\}$.
  B2 has the predecessors B1 and B5, and $\mathrm{OUT}[\mathrm{B5}]$ is
  still empty, so $\mathrm{IN}[\mathrm{B2}]=\{d_1,d_2\}$ and
  $\mathrm{OUT}[\mathrm{B2}]=\{d_3\}\cup(\{d_1,d_2\}-\{d_5\})=\{d_1,d_2,d_3\}$.
  B3 kills $d_2$ and adds $d_4$: $\{d_1,d_3,d_4\}$.  B4 kills $d_2$ and
  $d_3$ and adds $d_5,d_6$: $\{d_1,d_5,d_6\}$.  At the join B5 the union
  of the two is $\{d_1,d_3,d_4,d_5,d_6\}$, and B5 replaces $d_1$ by
  $d_7$: $\mathrm{OUT}[\mathrm{B5}]=\{d_3,d_4,d_5,d_6,d_7\}$, which also
  reaches EXIT.
\end{example}

\begin{example}[The second and third passes]
  \label{ex:l10-pass2}
  In pass~2, $\mathrm{OUT}[\mathrm{B5}]$ is no longer empty, and the
  back edge brings its definitions into B2:
  $\mathrm{IN}[\mathrm{B2}]=\{d_1,d_2\}\cup\{d_3,\dots,d_7\}$, all seven.
  B2 kills only $d_5$, so $\mathrm{OUT}[\mathrm{B2}]$ gains $d_4$, $d_6$
  and $d_7$, and so do the $\mathrm{IN}$ sets of B3 and B4.  Of these,
  B3 and B4 pass on only $d_7$, since each kills or regenerates the
  other two, and at B5 only $d_7$ is new, which B5 generates anyway:
  $\mathrm{OUT}[\mathrm{B5}]$ is unchanged.  Because B2, B3 and B4
  changed, a third pass is run; its $\mathrm{IN}[\mathrm{B2}]$ is the
  same as in pass~2, so no set changes, and the states of pass~2 are
  the fixed point.
\end{example}

\Needspace{26\baselineskip}
The solution answers questions an optimizer asks.  Both definitions of
\texttt{i}, $d_1$ and $d_7$, reach the use \texttt{x[i]} in B2, so
\texttt{i} cannot be replaced there by the constant~0; without the back
edge only $d_1$ would reach it.\margin{Only $d_3$ among the definitions
of \texttt{t} reaches B3 and B4: $d_5$ arrives at B2 through the back
edge but is killed there.  Each use of \texttt{t} in B3 and B4 thus
reads the value loaded in the same iteration.}  The initialisation
$d_2$ (\texttt{s = 0}) reaches B2 but not B5: both paths through the
loop body redefine \texttt{s}.  The definitions that reach EXIT,
$d_3$ to $d_7$, are the ones whose values the function can still
observe at its end.

\begin{keypoint}[Lesson summary]
  A compiler consists of a front end that builds an AST and an IR, an
  optimizer that rewrites the IR, and a back end that generates code.
  For CUDA, nvcc splits host from device code; the device code becomes
  PTX through cicc (or Clang and the LLVM NVPTX back end), PTX becomes
  SASS through ptxas, and a fat binary built by nvcc carries both, so that the driver
  can JIT-compile the PTX for GPUs that did not exist at build time.
  PTX is a virtual ISA with typed instructions of zero to four operands,
  unlimited virtual registers and guard predicates set by
  \texttt{setp}.  Optimisers work on three-address code divided into
  basic blocks, found with the three leader rules, and connected into a
  control-flow graph.  Global optimisations rest on data-flow analysis:
  transfer functions per block, a meet at joins, and an iterative
  algorithm that reaches a fixed point.  Reaching definitions is the
  forward analysis with
  $\mathrm{OUT}[B]=\mathrm{gen}_B\cup(\mathrm{IN}[B]-\mathrm{kill}_B)$
  and $\mathrm{IN}[B]$ the union over the predecessors, computed
  efficiently with bit vectors; it is fixed by its domain, direction,
  transfer function, meet, boundary condition and initial value, the
  six items that Lesson~11 fills in for liveness.
\end{keypoint}

\end{document}
